\documentclass[11pt]{article}

\usepackage[T1]{fontenc}
\usepackage[margin=1in]{geometry}
\usepackage{lmodern}
\usepackage{microtype}
\usepackage{amsmath,amssymb,mathtools}
\usepackage{amsthm}
\usepackage{aliascnt}
\usepackage{booktabs}
\usepackage{tabularx}
\usepackage[font=small,labelfont=bf]{caption}
\usepackage[section]{placeins}
\usepackage{float}
\usepackage{array}
\usepackage{xcolor}
\usepackage{enumitem}
\usepackage[hidelinks]{hyperref}
\usepackage[nameinlink,noabbrev]{cleveref}
\usepackage{url}
\usepackage[backend=biber,style=numeric,sorting=nyt,maxbibnames=99]{biblatex}
\addbibresource{references.bib}


\newtheorem{theorem}{Theorem}

\newaliascnt{lemma}{theorem}

\newtheorem{lemma}[lemma]{Lemma}
\aliascntresetthe{lemma}

\newaliascnt{proposition}{theorem}

\newtheorem{proposition}[proposition]{Proposition}
\aliascntresetthe{proposition}

\newaliascnt{corollary}{theorem}

\newtheorem{corollary}[corollary]{Corollary}
\aliascntresetthe{corollary}
\crefname{theorem}{Theorem}{Theorems}
\crefname{lemma}{Lemma}{Lemmas}
\crefname{proposition}{Proposition}{Propositions}
\crefname{corollary}{Corollary}{Corollaries}

\definecolor{deepblue}{RGB}{23,69,124}
\definecolor{lightblue}{RGB}{242,247,252}
\definecolor{darkred}{RGB}{145,32,32}
\ifdefined\ANONYMOUS

\newcommand{\pdfpaperauthor}{Anonymous}
\else

\newcommand{\pdfpaperauthor}{Justin Thaler, Kai Zhang, Scott Kominers, Quang Dao}
\fi
\hypersetup{
  colorlinks=true,
  linkcolor=deepblue,
  citecolor=deepblue,
  urlcolor=deepblue,
  pdftitle={A Symmetric-Square Quotient in Odd-Characteristic SNOVA and Conditional Forgery Cost Estimates},
  pdfauthor={\pdfpaperauthor},
  pdfsubject={Cryptanalysis of SNOVA Version 2.3},
  pdfkeywords={SNOVA, multivariate signatures, forgery attack, NIST PQC}
}
\setlist{leftmargin=*,topsep=3pt,itemsep=1.5pt,parsep=0pt,partopsep=0pt}


\newcommand{\F}{\mathbb F}

\newcommand{\Mat}{\operatorname{Mat}}

\newcommand{\Sym}{\operatorname{Sym}}

\newcommand{\rank}{\operatorname{rank}}

\newcommand{\Span}{\operatorname{span}}

\newcommand{\transpose}{\mathsf T}

\newcommand{\SNOVA}{\textsf{SNOVA}}

\newcommand{\LoneModelRange}{$2^{128.82}$--$2^{129.44}$}

\newcommand{\LthreeModelRange}{$2^{171.69}$--$2^{183.44}$}

\newcommand{\LfiveModelRange}{$2^{214.12}$--$2^{237.84}$}

\newcommand{\LoneAXNRange}{$2^{132.17}$--$2^{133.75}$}

\newcommand{\LthreeAXNRange}{$2^{175.04}$--$2^{187.74}$}

\newcommand{\LfiveAXNRange}{$2^{217.48}$--$2^{242.15}$}

\newcommand{\LoneNANDRange}{$2^{133.04}$--$2^{134.67}$}

\newcommand{\LthreeNANDRange}{$2^{175.91}$--$2^{188.67}$}

\newcommand{\LfiveNANDRange}{$2^{218.35}$--$2^{243.07}$}

\newcommand{\AXNNominalGapRange}{9.25--54.52}

\newcommand{\AXNExpectedAESGapRange}{8.13--52.86}

\newcommand{\AXNFullAESGapRange}{9.13--53.86}

\newcommand{\attacktitle}{A Symmetric-Square Quotient in
  Odd-Characteristic \texorpdfstring{\SNOVA}{SNOVA} and Conditional
  Forgery Cost Estimates}

\title{\attacktitle}
\ifdefined\ANONYMOUS
  \author{Anonymous submission}
\else
  \author{%
    Justin Thaler\thanks{a16z crypto research and Georgetown University} \and
    Kai Zhang\thanks{MIT} \and
    Scott Duke Kominers\thanks{Harvard University and a16z crypto research} \and
    Quang Dao\thanks{Carnegie Mellon University and LayerZero Labs}%
  }
\fi
\date{28 July 2026}

\begin{document}
\setlength{\emergencystretch}{2em}
\raggedbottom
\maketitle

\ifdefined\DISCLOSURE
\begin{center}
\fcolorbox{darkred}{white}{\parbox{0.91\textwidth}{\centering\small
\textbf{Confidential coordinated disclosure.} This draft is intended for the
\SNOVA{} submission team and NIST.\@ Please do not redistribute before the
coordinated public-disclosure date.}}
\end{center}
\fi

\begin{abstract}
At its algebraic core, a \SNOVA{} forgery requires solving a public system of
multivariate quadratic (MQ) equations; the resulting assignment must also pass
the verifier's rejection checks.  We show that the public $q=19$ design
represents many of its quadratic features twice.  A feature is labeled by an \emph{ordered} pair of
powers from the public algebra $A=\F_q[S]$, but symmetry of $S$ and of the
public base matrices makes the polynomial unchanged when the two labels are reversed.
Thus, if $d=\dim_{\F_q}A$, the $d^2$ labeled power-pair features of each
underlying public form factor through at most $\binom{d+1}{2}$ unordered
features.  In the public parameter sets this
reduces $16$ labels to $10$ at block size $\ell=4$, and $4$ labels to $3$
at block size $\ell=2$.  The identification occurs before the public outer map and holds for
every symmetric public key; it is not a weak-key event.

For each of the nine public $q=19$ parameter sets in \SNOVA{}
Version~2.3, we give an attack construction that, when its explicit public
preflights pass, restricts the verifier to an attacker-chosen column relation
and reduces the resulting system to a smaller root system while retaining all
$M$ public output equations.  For $\ell=4$, the chosen column relation
preserves the four-block structure used by special XL.\@  For $\ell=2$, a
fixed-skew relation ensures that every root recovered from the residual system
reconstructs to a verifier root that passes the block-symmetry rejection test.
The quotient, affine reduction, rejection
identities, probability bounds, and final verification wrapper are exact;
uniform-target and independent-retry statements are exact within the
explicitly stated random-oracle experiment.  Under unproved global-rank and
production-solver hypotheses, the modeled expected operation counts at NIST
Levels~I, III, and V are, respectively, \LoneModelRange,
\LthreeModelRange, and \LfiveModelRange.  Replacing the field arithmetic
charged by those models with author-reported two-input AND/XOR/XNOR circuits,
validated as described in \Cref{app:field-circuits}, gives
AXN arithmetic-schedule counts \LoneAXNRange, \LthreeAXNRange, and
\LfiveAXNRange.  These are neither wall-clock forgeries nor complete attack
circuits: production solving behavior, control logic, matrix construction,
memory, communication, and data movement are not priced.
\end{abstract}

\section{Introduction}\label{sec:intro}

\paragraph{The cryptanalytic task.}
A multivariate signature scheme publishes a polynomial map of total degree
at most two,
\[
  \mathcal P:\F_q^{N}\longrightarrow\F_q^{M}.
\]
For a message with hash-derived target $t\in\F_q^M$, the algebraic core of a
forgery is to find $x$ satisfying $\mathcal P(x)=t$; any additional format or
rejection checks must also pass.  This is an MQ problem: $M$ equations of
degree at most two in the coordinates of $x$.  Linear and constant terms are
allowed; throughout the paper, ``quadratic'' refers to the maximum degree.
The apparent equation and variable counts are only a first approximation to
security.  Public algebraic structure can make different-looking quadratic
features identical, can expose linear equations, or can split the remaining
variables into blocks that a specialized solver can exploit.

\SNOVA{} is one of the nine digital-signature candidates that NIST advanced to
the third round of its additional post-quantum signature process in
May~2026~\cite{NISTThirdRound2026}.  It stores the signature variables in small
matrices and builds its public quadratics from pairs of powers of a public
matrix $S$.  Beullens showed that an attacker can impose an affine relation
among the signature columns and thereby reduce forgery to a smaller
underdetermined MQ system \cite{Beullens2024SNOVA}.  In that analysis he also
observed that reversed bilinear terms can be combined for a symmetric MAYO
base map, while the then-analyzed nonsymmetric \SNOVA{} matrices required
ordered terms.  The symmetrized public matrices of Version~2.3 change that
premise.  We formalize the resulting deterministic quotient inside the
current public quadratic features and integrate it with an exact affine
reduction of all $M$ equations after the attacker imposes a column relation.

\paragraph{The hidden duplication.}
Assume $S^\transpose=S$, let $P_i$ be a public symmetric base matrix, put
$A=\F_q[S]$, and write $\Sigma_\xi=I_n\otimes\xi$ for $\xi\in A$.  The feature
with left label $\xi$ and right label $\eta$ is
\[
  q_{i,\xi,\eta}(u)
  =u^\transpose\Sigma_\xi P_i\Sigma_\eta u.
\]
Because the matrices are symmetric and the result is a scalar,
\begin{equation}\label{eq:intro-swap}
  q_{i,\xi,\eta}(u)=q_{i,\eta,\xi}(u).
\end{equation}
The verifier begins from ordered labels, but the polynomial depends only on
the corresponding unordered pair.  We call the operation of identifying
these duplicate labels the \emph{symmetry quotient}.

For a two-dimensional toy algebra with basis $\{1,S\}$, the four labels
\[
 (1,1),\qquad(1,S),\qquad(S,1),\qquad(S,S)
\]
produce only three polynomials because the middle two are equal.  More
generally, if $\dim_{\F_q}A=d$ and $m_1$ is the number of public base-form
groups, then at most
\begin{equation}\label{eq:intro-K}
  K=m_1\binom{d+1}{2}
\end{equation}
power-pair quotient coordinates remain across the $m_1$ base forms, rather
than $m_1d^2$.  Thus each $\ell=4$ base form drops from $16$ ordered labels to
$10$ unordered labels; each $\ell=2$ base form drops from $4$ to $3$.  The
loss occurs before the public outer linear map---the final public mixing of
these features into verifier outputs---so later randomization cannot recreate
the discarded directions.  It holds for every symmetric public key
and every attacker-chosen column relation.

\begin{center}
\fcolorbox{deepblue}{lightblue}{\parbox{0.94\textwidth}{\small
\textbf{Core idea in plain language.}  The verifier assigns two names to the
same quadratic polynomial whenever it reverses a pair of power labels.  The
attack first removes those duplicate features, then uses ordinary linear
algebra to retain all remaining verifier equations, and finally solves only
the smaller quadratic core.}}
\end{center}

\paragraph{The attack plan.}
The proof is easiest to follow as a sequence of transformations.  The
word \emph{lift} below means the attacker-chosen affine rule that turns one
column vector into a full signature candidate.  Concretely, if
$u_j\in\F_q^{N_0}$ is the stacked $j$th column of the signature, the attacker
imposes
\[
  u_j=\Sigma_{R_j}u+v_j\qquad(0\le j<r),
\]
where $u\in\F_q^{N_0}$ is one common column variable and the public
multipliers $R_j\in A$ and offsets $v_j\in\F_q^{N_0}$ are chosen before the
MQ solve.  Thus the search moves from $rN_0$ column coordinates to one
$N_0$-coordinate variable plus fixed offsets; \Cref{sec:affine-columns}
defines the relation and its feature maps formally.

The equality \eqref{eq:intro-swap} is only the first job.  The construction
also has to preserve all public outputs, ensure that targets have roots, and
make the residual equations compatible with a practical solver.  It has four
steps:
\begin{enumerate}[label=\textbf{\arabic*.},itemsep=3pt]
\item \textbf{Constrain the signature columns.}  Every signature column is
      written as an affine function of one common vector.  Substitution
      separates the verifier into quotient quadratic features, offset-linear
      terms, and constants.
\item \textbf{Eliminate the affine part without dropping equations.}  Output
      directions outside the quadratic image become linear constraints.  For
      each target satisfying those affine consistency conditions, Gaussian
      elimination leaves a smaller MQ system with exactly the same solutions
      as the constrained verifier.
\item \textbf{Choose a solver-friendly search space.}  For $\ell=4$, a
      structured choice of offsets leaves a large $A$-linear kernel.  Over a
      splitting field this kernel becomes four variable blocks, preserving the
      multidegrees used by special XL.\@  For $\ell=2$, a fixed-skew relation
      handles the verifier's block-symmetry rejection test directly.  A
      zero-offset alternative is analyzed separately, but is not used in the
      headline table because a dense accepted-root population does not by
      itself determine which root an MQ solver returns.
\item \textbf{Retry and verify.}  Exact moment bounds control the
      probabilities of root existence and uniqueness for random target--slice
      pairs.  A candidate recovered by the MQ solver is transformed back and
      checked by the original verifier.
\end{enumerate}

After the attacker-chosen column relation is imposed, every transformation
through the residual MQ system is exact and uses only public algebra and
linear algebra.  The solver model enters only when we estimate the cost of
solving that residual system.  The final public-verifier check makes the
wrapper zero-error: a trial may be discarded, but an invalid signature is
never returned.  It is a finite-expected-time Las Vegas attack only under the
stated independent-trial, global-rank, and solver-success hypotheses.

\paragraph{A complete dimension walk-through.}
The Level-I parameter set $(v,o,\ell,r)=(28,5,4,4)$ illustrates where the
reduction comes from: it has 28 vinegar blocks, 5 oil blocks, and blocks with
4 rows and 4 columns.  Here $n=v+o=33$, one stacked signature column has
$N_0=n\ell=132$ coordinates, the verifier has $M=or\ell=80$ scalar outputs,
and there are $m_1=5$ public base-form groups.

\begin{table}[H]
\centering
\caption{The dimensions in the $(28,5,4,4)$ attack, conditional on the stated
public rank preflights.  This table is only a worked example; the general
formulas appear in later sections.}
\label{tab:running-example}
\small
\renewcommand{\arraystretch}{1.12}
\begin{tabularx}{0.96\textwidth}{@{}>{\raggedright\arraybackslash}p{0.28\textwidth}>{\centering\arraybackslash}p{0.18\textwidth}X@{}}
\toprule
Stage & Dimension & Meaning\\
\midrule
Raw public feature vector & $N_{\rm raw}=1{,}280$ & $m_1\ell^2r^2$ labeled slots before
the column relation\\
After the common-column relation & $80$ ordered quadratics & one feature for
each of $m_1\ell^2=5\cdot16$ ordered power pairs\\
After the symmetry quotient & $K=50$ quotient coordinates &
$m_1\binom{\ell+1}{2}=5\cdot10$ unordered power pairs; the public preflight
checks $\rank Q_R=50$\\
Affine output directions & $M-K=30$ & linear equations that determine an
affine search space\\
Full residual domain & $102$ coordinates & $N_0-(M-K)=132-30$ after affine
elimination\\
Stable kernel & $52$ base-field dimensions & an $A$-linear space of
$A$-dimension $13$ on which offset-linear motion vanishes\\
Selected special-XL slice & $h=40$ & four blocks of $s=10$ variables; hence
$\tau=K-h=10$, with restricted polynomial rank
$\rho_{\rm poly}(L)=50$ checked on the selected slice\\
\bottomrule
\end{tabularx}
\end{table}

The final slice deliberately has fewer variables than equations.  Under full
cross-polar rank, a random target--slice trial then contains a unique
base-field root with probability essentially $19^{-10}$.  Paying roughly
$19^{10}$ retries is worthwhile because reducing the Macaulay matrix from the
full residual domain to four blocks of ten variables saves much more.  This is
why the paper analyzes both the algebraic reduction and the root-retry
probability, rather than quoting an MQ dimension alone.

\paragraph{Headline estimates.}
\Cref{tab:main} reports modeled costs and the differences between them.
``Model'' is the base-2 expected-cost exponent in the \emph{pinned} attack
model, meaning that we fix the cited estimator, its version, and its conventions
exactly as specified in
\Cref{app:cost-certificates,app:generic-mq-certificate}.  It starts from the cited per-solve MQ
operation formula and, when the solve is repeated, includes
the explicit target--slice retry factor.  Target generation performed before a
single solve is accounted for separately.  ``AXN arith.\ schedule'' changes
only the unit for the arithmetic charged by that model, replacing its coarse
field-operation factor with signed two-input AND/XOR/XNOR (AXN) circuits.  ``Internal estimator gain'' is the
reduction in the Model exponent
relative to a pinned generic-MQ baseline on the same symmetry quotient.  It is
a solver improvement only in the structured rows, where it comes from the
stable four-block slice and special XL.\@  ``Nominal arithmetic gap'' is the
nominal NIST reference minus the AXN arithmetic-schedule exponent.  It
is room for unpriced overhead in that one metric, not an unconditional
security margin.

\begin{table}[H]
\centering
\caption{Headline conditional estimates for the public Version~2.3 parameter
sets.  The tuple is $(v,o,\ell,r)$; its meaning is given in
\Cref{sec:parameters}.  Model and AXN are base-2 cost exponents; internal
estimator gain and nominal arithmetic gap are differences of exponents,
measured in bits.
V2.3 reproduces the historical ``log \#gates'' values reported with the
Version~2.3 materials~\cite{SNOVA23}; it is contextual rather than a
same-dimension recosting.  The superscript
$\mathsf M$ marks quantities conditional on the global-rank and
production-solver hypotheses stated in the text; both derived-difference
columns inherit the same conditions.  Every row is additionally conditional
on the attacked public key admitting the required quotient-rank, affine-rank,
slice (where applicable), and rejection preflights for that mode; the
manuscript does not prove that a passing choice exists for every valid public
key.}
\label{tab:main}
\footnotesize
\setlength{\tabcolsep}{2.3pt}
\renewcommand{\arraystretch}{1.08}
\begin{tabular}{@{}c l l r r r r r@{}}
\toprule
Level & Parameters & Mode & V2.3 & Model$^{\mathsf M}$ &
\shortstack{AXN arith.\\schedule$^{\mathsf M}$} &
\shortstack{internal estimator\\gain} & \shortstack{Nominal arith.\\gap}\\
\midrule
I   & $(28,5,4,4)$  & structured XL & 194 & 128.82 & 132.17 &  7.22 & 10.83\\
I   & $(48,16,2,2)$ & fixed-skew MQ & 164 & 129.44 & 133.75 &  0.00 &  9.25\\
I   & $(28,4,4,5)$  & structured XL & 194 & 128.82 & 132.17 &  7.22 & 10.83\\
\midrule
III & $(40,7,4,4)$  & structured XL & 264 & 171.69 & 175.04 &  8.91 & 31.96\\
III & $(72,24,2,2)$ & fixed-skew MQ & 234 & 183.44 & 187.74 &  0.00 & 19.26\\
III & $(38,5,4,5)$  & structured XL & 239 & 171.69 & 175.04 &  8.91 & 31.96\\
\midrule
V   & $(50,9,4,4)$  & structured XL & 333 & 214.12 & 217.48 & 10.56 & 54.52\\
V   & $(96,32,2,2)$ & fixed-skew MQ & 303 & 237.84 & 242.15 &  0.00 & 29.85\\
V   & $(52,6,4,6)$  & structured XL & 333 & 214.12 & 217.48 & 10.56 & 54.52\\
\bottomrule
\end{tabular}
\end{table}

For example, the first row's AXN entry $132.17$ means about $2^{132.17}$
charged two-input arithmetic gates under the special-XL model and its
unique-root retry factor.  It does \emph{not} include a complete implementation
of sparse matrix construction, pivot control, memory traffic, or data
movement.

\paragraph{Conditions behind the table.}
The structured-XL rows require two separate production assumptions: every
nonzero output direction must have full cross-polar rank on the chosen slice,
and a planted trial conditioned on having a unique $\F_q$-root must produce
the predicted one-dimensional special-XL Macaulay right kernel.  The
three fixed-skew rows use the sufficient all-target root conditions
$r_*\ge96$, $r_*\ge144$, and $r_*\ge192$.  Each $r_*$ condition is global: it
concerns every nonzero output combination, not only the sampled directions in
the reported finite campaigns.  Lower ranks are handled by explicit retry bounds in
\Cref{sec:slices}.  The zero-offset Level-V alternative is summarized in
\Cref{sec:l2-costs} and analyzed in \Cref{app:zero-offset}; it is not a
headline estimate because \Cref{prop:l2-rejection} counts accepted roots but
does not specify the output distribution of the priced MQ solver.

\begin{center}
\fcolorbox{deepblue}{lightblue}{\parbox{0.94\textwidth}{\small
\textbf{Five levels of claim.}\par\smallskip
\begin{tabularx}{\linewidth}{@{}>{\bfseries}p{0.24\linewidth}X@{}}
Exact mathematics & Symmetry quotient, affine elimination, stable-kernel
construction, fixed-skew identity, rejection identities, probability
inequalities, and final-verifier wrapper.\\[2pt]
Random-oracle exactness & Rejection conversion gives uniform residual targets
and fresh trials are independent, conditional on the explicit random-oracle
experiment and duplicate-query check in \Cref{sec:target-sampling}.\\[2pt]
Public preflight & Deterministic rank and rejection checks for the chosen key,
relation, offsets, and slice.  A failing choice is discarded before an
expensive MQ solve.\\[2pt]
Finite evidence & Author-reported, nonexhaustive rank campaigns and smaller
Macaulay experiments.  They support, but do not certify, the global production
assumptions and require the release contract in \Cref{app:artifact} for
independent reproduction.\\[2pt]
Cost model & Published generic-MQ or special-XL arithmetic schedules with an
explicit retry factor and a validated circuit cost for the charged field
primitive.  A complete solver implementation is not priced.
\end{tabularx}}}
\end{center}

The official Version~2.3 values are numerically $30.25$--$115.52$ bits above
the AXN arithmetic column, but the two sets of figures do not come from a
shared complete-circuit methodology.  A cleaner internal comparison recosts a
historical ordered-label affine-column estimator in the same AXN basis; the
resulting per-row differences are $38.02$--$129.66$ bits
(\Cref{tab:ordered-baseline}).  The unreduced affine-column equations are an
executable system, but \Cref{thm:symmetric-square} proves that their reversed
power labels do not represent independent quadratic features on a symmetric
key.  What is counterfactual is feeding the ordered-label count into a generic
estimator as though those labels described a generic quadratic space.  The
comparison is included only to isolate the numerical effect of recognizing
the quotient and then choosing the selected post-quotient attack mode.  It is
not a row-by-row speedup over the best prior Version~2.3 attack after all
public duplicate-column preprocessing.

The rectangular $(38,5,4,5)$ row needs an additional convention warning.  The
pinned reference implementation uses
$m_1=\lceil or/\ell\rceil=7$.  The authors report that the accompanying
historical analysis script instead uses the floor value $6$; an archival
release should identify that script's exact path and revision.  Our dimensions
use the implementation ceiling, so the historical V2.3 number in the table is
especially unsuitable as a same-instance difference.  The immutable
implementation revision used here is the one cited in \cite{SNOVA23}.

NIST describes its category references across the security metrics it regards
as relevant, rather than as a gate-count-only test~\cite{NISTPQCCriteria}.  We
therefore use $143/207/272$ only as nominal arithmetic reference points.
Establishing an end-to-end category break would require stronger production
rank evidence or measurements and a substantially more complete
implementation-level cost account.

\paragraph{How to read the paper.}
A reader new to MQ can read \Cref{sec:mq-background} first, then focus on the
plain-language paragraphs before and after each theorem.  The exact algebraic
reduction is in \Cref{sec:prelim,sec:symmetry,sec:structured-lift}.  Root existence and
retry probabilities are separated in \Cref{sec:slices}; the solver and gate
models are isolated in \Cref{sec:complexity}; and
\Cref{sec:evidence} states exactly which assumptions are tested rather than
proved.  The appendices contain the deferred proofs, the unselected
zero-offset analysis, implementation and circuit details, dimension formulas,
cost-search certificates, and the reproducibility contract.

\section{MQ Background}\label{sec:mq-background}

This section supplies the minimum background needed for the attack.  Readers
already comfortable with underdetermined MQ, polar forms, Macaulay matrices,
and multigraded XL may skip to \Cref{sec:prelim}.

\subsection{MQ systems, roots, and affine slices}\label{sec:mq-primer}

Write $\F_q$ for the finite field with $q$ elements.  An MQ system over
$\F_q$ is a vector-valued polynomial map of total degree at most two,
\[
  g=(g_1,\ldots,g_K):\F_q^N\longrightarrow\F_q^K,
\]
together with a target $z\in\F_q^K$.  We continue to call $g$ quadratic when
some coordinates also contain linear or constant terms.  A \emph{root} is an
$x$ with $g(x)=z$.
The set
\[
  g^{-1}(z)=\{x:g(x)=z\}
\]
is the \emph{fiber} above $z$: it is the set of all candidate assignments
that produce the target.  After the affine reconstruction described later,
those assignments become full signature candidates.  Every
production-parameter \SNOVA{} calculation uses $q=19$; the fully enumerated
theorem checks in \Cref{sec:evidence} also use $\F_3$ and $\F_5$.

As a toy example, two equations such as
\[
 x_1x_2+x_3^2=z_1,
 \qquad
 x_1^2+x_2x_3+x_1=z_2
\]
form an MQ system in three variables.  Having more variables than equations
makes the system \emph{underdetermined}, but does not make it linear or
necessarily easy.  Modern MQ attacks often spend most of their time on very
large linear-algebra matrices built from the quadratic equations.

An \emph{affine space} is a translate $a+V=\{a+v:v\in V\}$ of a linear
subspace.  An \emph{affine slice} $a+L\subseteq a+V$ restricts the search to a
smaller direction space $L$.  If $h=\dim L$, choosing a basis of $L$ turns the
slice into an MQ system in $h$ coordinates.  Smaller $h$ makes the algebraic
solve cheaper, but it also lowers the chance that the slice contains a root;
\Cref{sec:slices} quantifies this tradeoff.

Three logically distinct events recur throughout the paper:
\begin{enumerate}[label=(\roman*),itemsep=2pt]
\item the chosen target and slice contain at least one $\F_q$-rational root;
\item the slice contains exactly one such root; and
\item the particular Macaulay matrix used by special XL has a one-dimensional
      right kernel from which row reduction reveals that root.
\end{enumerate}
The first two are properties of the quadratic map.  The third is a statement
about a specific solver representation.  A proof of root existence is not, by
itself, a proof that one solver invocation recovers the root.

\subsection{Rank, polar forms, and collisions}\label{sec:rank-primer}

For a scalar quadratic polynomial $f$, define its polar form by
\[
  \beta_f(x,y)=f(x+y)-f(x)-f(y)+f(0).
\]
The polar form is bilinear in every characteristic.  In characteristic two it
annihilates square terms; the Gauss-sum fiber bounds used later are stated only
for odd $q$.  It is the finite-field analogue of the part of a directional
derivative that comes from the quadratic terms.  Its matrix rank measures how
many independent directions the quadratic part couples.

For a vector-valued map $g=(g_1,\ldots,g_K)$, every nonzero coefficient vector
$b\in\F_q^K$ selects a scalar combination $b\cdot g$.  The minimum polar rank
must be taken over \emph{all} such output directions.  Multiplying $b$ by a nonzero scalar does not change the rank, so these
vectors can be grouped into projective directions: one-dimensional lines
through the origin.

Why is rank useful?  Finite-field Fourier, or character-sum, identities imply
that a scalar quadratic with high polar rank is well balanced.  Applied to all nonzero combinations
$b\cdot g$, this gives rigorous bounds on every target fiber; when the ranks
are high enough relative to $K$, its size is close to the random-map prediction
$q^{N-K}$.  On a slice, a related
\emph{cross-polar rank} measures how strongly a displacement inside the slice
interacts with the full residual domain.  High cross-polar rank suppresses
pairs of distinct slice points that map to the same target.  Thus:
\begin{center}
\begin{tabularx}{0.91\textwidth}{@{}>{\bfseries}p{0.28\textwidth}X@{}}
Polar rank & controls whether full-space target fibers are nonempty and close
to their expected size.\\
Cross-polar rank & controls collisions inside an affine slice and therefore
unique-root retry probabilities.\\
Macaulay kernel & its dimension controls whether the chosen special-XL linear
algebra actually exposes the root.
\end{tabularx}
\end{center}
High rank is not invoked as a vague ``randomness'' heuristic in the structural
proofs: \Cref{sec:slices} gives exact first- and second-moment formulas.  The
production issue is whether the required rank lower bound holds in every
nonzero output direction.  Public sampling gives evidence, but only a proof or
exhaustive enumeration certifies the global minimum.

\subsection{XL and Macaulay matrices in one page}\label{sec:xl-primer}

XL (``eXtended Linearization'') turns a nonlinear system into a larger linear
system.  Starting from quadratic equations $f_i=0$, it multiplies each $f_i$
by selected monomials.  For example,
\[
  f=x_1x_2+x_3^2+x_1+1
  \quad\Longrightarrow\quad
  x_1f=x_1^2x_2+x_1x_3^2+x_1^2+x_1.
\]
Each resulting polynomial becomes a row.  Each monomial---for example
$x_1^2x_2$ or $x_1x_3^2$---becomes a column.  The coefficient matrix is a
\emph{Macaulay matrix}.  Gaussian elimination on this matrix can reveal enough
relations among the monomials to recover a root.  Treating monomials as columns
is only a linear-algebra device: the monomials are not independent variables in
the original problem, and the solver must eventually recover coordinates that
satisfy all of their multiplicative relations.

The cost comes from the size and rank of this matrix, not from evaluating the
original equations once.  Generic XL tracks only total degree.\@  \emph{Special
XL} exploits a variable-block structure.  A monomial then has a vector of
degrees, one entry per block, called its \emph{multidegree}.  Building only
selected multidegrees can make the Macaulay matrix much smaller.

Our $\ell=4$ reduction produces four blocks after a change to eigenblock
coordinates.  A quadratic using one block twice has degree pattern
$(2,0,0,0)$ up to permutation; a bilinear equation using two blocks has pattern
$(1,1,0,0)$.  The stable-lift construction in
\Cref{sec:structured-lift} exists precisely to preserve these patterns after
affine elimination.  A \emph{homogenizer} is one extra variable per block that
allows the constant and linear terms to be written with the same
multidegree; it is set to one when the root is read out.

If $x$ is a root, the vector of all selected monomial values
\[
  \nu(x)=(1,x_1,\ldots,x_N,x_1^2,x_1x_2,\ldots)
\]
lies in the right kernel of the dehomogenized Macaulay matrix.  The right
kernel is the space of column vectors $w$ with $\mathcal M w=0$.  If it is
one-dimensional, normalizing its constant coordinate to one reveals the
degree-one coordinates and hence the root.  A unique base-field root does
not automatically imply a one-dimensional Macaulay kernel: extension-field
roots or algebraic relations among rows can create extra kernel vectors.  This
is the separate production solver assumption in the structured attacks.

\section{Relation to Prior Attacks}\label{sec:related}

Beullens rewrote the \SNOVA{} verifier in ordered bilinear power-pair
coordinates and derived affine-column forgery attacks from rank defects in the
resulting outer matrix~\cite{Beullens2024SNOVA}.  Our quotient acts one layer
earlier: public-matrix symmetry identifies the power-pair polynomials
\emph{before} the outer map sees them.  The underlying reversed-term
observation is not unprecedented: Beullens explicitly notes that the
symmetric MAYO base map permits combining the two orders, and contrasts that
case with the nonsymmetric \SNOVA{} matrices then under analysis.  Our narrow
novelty claim is the application to the symmetrized Version~2.3 public
matrices and its integration with exact affine elimination, stable slices,
rejection handling, and retry bounds.

Cabarcas et al.\ exploit stability under the public matrix algebra with a
multigraded XL algorithm~\cite{Cabarcas2025SNOVA}.  Furue et al.\ expose related
multihomogeneous systems through matrix transformations~\cite{FurueEtAl2026}.
Our attack starts from all $M$ equations of the verifier after imposing a
public column relation, uses exact affine elimination, and then
restricts to an invariant kernel on which the affine offsets preserve the
block grading needed by special XL.\@

Two more recent formulation papers are especially relevant to the boundary of
this contribution.  Li, Guo, and Ding place the unbalanced-oil-and-vinegar
(UOV)-like schemes QR-UOV, MAYO, and \SNOVA{} in a common tensor-product
framework and use it primarily for scheme comparison and
design~\cite{LiGuoDing2025Unified}.  Ding et al.\ give ring-equation,
whipping, and tensor formulations of \SNOVA{} and propose a more flexible
reformulated design, including odd-characteristic parameters
\cite{DingEtAl2026Reformulating}.  The point claimed here is not priority for
tensor language or for symmetry in UOV-like schemes.  It is the explicit
kernel of the ordered label map for the symmetric Version~2.3 verifier and the
consequent public direct-forgery reduction.  In particular, the large
ordered-label comparison in \Cref{tab:ordered-baseline} is not a
state-of-the-art comparison with these works.

\begin{table}[tbp]
\centering
\caption{Scope of the closest formulation and attack lines.  ``Ordered'' and
``symmetric-square'' describe the public power-pair feature representation,
not the full internal tensor language of each work.}
\label{tab:related-scope}
\scriptsize
\setlength{\tabcolsep}{3.8pt}
\begin{tabularx}{0.97\textwidth}{@{}
>{\raggedright\arraybackslash}p{0.18\textwidth}
>{\raggedright\arraybackslash}p{0.19\textwidth}
>{\raggedright\arraybackslash}p{0.18\textwidth}
>{\raggedright\arraybackslash}X@{}}
\toprule
Work & Public-matrix premise & Primary goal & Relation to this paper\\
\midrule
Beullens~\cite{Beullens2024SNOVA} & Nonsymmetric \SNOVA{} base matrices in the
analyzed design & Direct forgery by affine columns & Uses ordered terms for
\SNOVA{} and already identifies their collapse for symmetric MAYO.\\
Cabarcas et al.~\cite{Cabarcas2025SNOVA} & Algebra-stable \SNOVA{} structure &
Multigraded forgery solving & Supplies the special-XL model used after our
quotient and stable-slice construction.\\
Li--Guo--Ding~\cite{LiGuoDing2025Unified} & Tensor-product UOV-like schemes &
Unified formulation and scheme design & Establishes nearby tensor language;
our claim is the concrete public quotient and direct-forgery reduction.\\
Ding et al.~\cite{DingEtAl2026Reformulating} & Ring, whipping, and tensor
forms of \SNOVA{} & Reformulation and parameter design & Motivates a
comparison with redesigned odd-characteristic systems; it is not used as a
prior attack baseline here.\\
This paper & Symmetric Version~2.3 base matrices & Exact quotient and
conditional direct-forgery costs & Factors ordered labels through
$\Sym^2(A)$ before the outer map and retains all $M$ equations of the
substituted verifier.\\
\bottomrule
\end{tabularx}
\end{table}

The last two formulation papers are not assigned residual dimensions or
solver exponents in our tables because we do not import an attack instance
from them.  Conversely, we do not claim that their tensor descriptions omit
every equivalent symmetry identity.  The concrete comparison made here is
only this: the present attack fixes the symmetric, odd-characteristic
Version~2.3 verifier, proves the quotient before the outer map, retains all
outputs, records its residual dimensions, and then applies the specified
generic-MQ or special-XL estimator.

The wedge attack breaks the submitted characteristic-two parameters through
the block-ring structure~\cite{BrosEtAl2026SNOVAWedge}; later work studies the
wedge map further~\cite{TsengEtAl2026Wedge}.  Other analyses reinterpret
scalar-expanded \SNOVA{} as UOV, or lift it to extension-field UOV, for key
recovery~%
\cite{IkematsuAkiyama2024,LiDing2024,NakamuraEtAl2025}.  Symmetric and
exterior algebras also appear in recent UOV key-recovery attacks~%
\cite{JinEtAl2026,Ran2026}.  Our attack is a direct forgery attack: it does not
recover the hidden oil subspace that constitutes the UOV-style secret key.

For the generic $\ell=2$ cores, we reproduce Hashimoto's optimization,
including the published boundary value $a=1$, where $a$ is the optimizer's
square-subsystem parameter~%
\cite{Hashimoto2023,Beullens2024SNOVA}.  When this leaves the elementary system
$\operatorname{MQ}(q,1,1)$, the pinned convention solves it by exhaustive evaluation and
assigns one primitive-unit charge per field element tested.  May, Ostuzzi, and Ressler's \emph{Just
Guess} algorithm and the generalized variable partitioning of Asanuma et al.\ are more recent approaches to underdetermined MQ~%
\cite{MayOstuzziRessler2026,AsanumaEtAl2026}.  Evaluating the published
classical Just-Guess formulas on our residual dimensions gives larger estimates
than the selected Beullens--Hashimoto modes.  We have not implemented the newer
generalized partition search.  The ``generic MQ'' numbers below are therefore
pinned, reproducible baselines rather than a claim that every possible generic
solver has been exhausted.

Solver engineering is also moving quickly.  Sakata and Takagi use
Hilbert-series information to reduce the matrices built by an $F_4$
Gr\"obner-basis solver and report a new MQ-challenge
record~\cite{SakataTakagi2026}.  That work is an
implementation and challenge result rather than a plug-in cost formula for the
underdetermined residual systems here, so we do not substitute it into
\Cref{tab:main}.  It reinforces the need to distinguish exact algebraic
reductions, solver-operation models, and complete implementations.


\section{The Affine-Column Framework}\label{sec:prelim}

This section rewrites the public verifier after an attacker-imposed affine
relation among signature columns.  The result has three visibly separate
parts: quotient quadratic features, linear terms created by the offsets, and a
constant term.  \Cref{prop:affine-reduction} then eliminates the affine part
exactly.

Three vector spaces appear and should not be confused.  The raw public feature
vector has length $N_{\rm raw}$; the verifier outputs $M$ scalar equations; and the
common column variable $u$ has length $N_0$.  The matrices below simply move
between these three spaces.

\subsection{Notation and dimensions}\label{sec:notation}

\begin{table}[H]
\centering
\caption{Recurring notation.  Generic primer notation such as $N$ is local;
$N_0$ is the specific common-column dimension in the \SNOVA{} reduction.}
\label{tab:notation}
\small
\renewcommand{\arraystretch}{1.12}
\begin{tabularx}{0.96\textwidth}{@{}>{\raggedright\arraybackslash}p{0.23\textwidth}X@{}}
\toprule
Symbol & Meaning\\
\midrule
$(v,o,\ell,r)$ & Public parameter tuple: $v$ vinegar blocks, $o$ oil blocks,
block height $\ell$, and $r$ columns per signature block.\\
$n=v+o$ & Number of matrix-valued signature blocks.\\
$M=or\ell$ & Number of scalar public output equations.\\
$N$ & Generic domain dimension for an MQ map $g:\F_q^N\to\F_q^K$; its
meaning is local to the section in which the map is defined.\\
$N_0=n\ell$ & Length of one stacked signature column after the common-column
relation.\\
$N_{\rm raw}=m_1\ell^2r^2$ & Number of labeled verifier feature slots before imposing
the relation.\\
$A=\F_q[S]$ & Public commutative matrix algebra generated by $S$; in the
Version~2.3 rows, $\dim_{\F_q}A=\ell$.\\
$m_1$ & Number of public base-form groups in the pinned implementation.\\
$K=m_1\binom{\ell+1}{2}$ & Number of quotient-feature coordinates supplied by
the symmetric-square representation; their polynomial span may have smaller
dimension until a rank preflight is performed.\\
$Q_R$, $\rho=\rank_{\F_q}Q_R$ & Linear outer map from the $K$ quotient
coordinates to public output directions, and the dimension of its image.
This is not the matrix rank of a scalar quadratic form.\\
$\rho_{\rm poly}(L)$ & Dimension of the span of the homogeneous quotient
coordinate polynomials after restriction to a slice $L$; this is the
``restricted feature rank.''\\
$C_{R,\mathbf v}$ & Linear constraint matrix obtained by projecting away the quadratic
image.\\
$H_{R,\mathbf v}$ & Stable $A$-linear kernel used to choose structured solver slices.\\
$L$, $h$ & Slice direction space and its base-field dimension.\\
$\tau=K-h$ & Slice deficit.  It enters the exact full-rank retry formula;
for the selected $\tau=10,14$ slices, the retry count is essentially $q^\tau$.\\
$r(b)$, $r_*$ & Matrix rank of the polar bilinear form of the scalar
combination $b\cdot g$, and its minimum over $b\ne0$.\\
$e_L(b)$ & Rank of the cross-polar linear map in output direction $b$
relative to the slice $L$.\\
$\mu$, $\bar\eta$ & Expected roots on a random target--slice pair and an upper
bound on its collision mass.\\
\bottomrule
\end{tabularx}
\end{table}

\subsection{Parameters and public forms}\label{sec:parameters}

The public parameter tuple is $(v,o,\ell,r)$.  There are $v$ vinegar blocks and
$o$ oil blocks, for a total of
\[
  n=v+o
\]
matrix-valued signature blocks.  The names ``vinegar'' and ``oil'' come from
SNOVA's UOV ancestry; in this public forgery attack they serve only to specify
dimensions, and the attacker never learns or reconstructs the secret partition.
Each block has $\ell$ rows and $r$ columns.
Stacking the $j$th column of all $n$ blocks gives a vector of length
$N_0=n\ell$; the verifier has $M=or\ell$ scalar outputs.  The pinned
Version~2.3 implementation defines
\begin{equation}\label{eq:parameters-basic}
 m_1=\left\lceil\frac{or}{\ell}\right\rceil,
 \qquad M=or\ell,
 \qquad N_0=n\ell.
\end{equation}
Here $m_1$ is the number of public base-form groups used to assemble the
outputs.  The ceiling is used by the reference implementation; a public
analysis script uses a floor and therefore disagrees when $\ell\nmid or$.  All
dimensions in this paper follow the pinned implementation~\cite{SNOVA23}.

Let $S\in\Mat_\ell(\F_q)$ be the public symmetric matrix with irreducible
characteristic polynomial.  Its powers generate the commutative matrix
algebra
\[
 A=\F_q[S],\qquad \Sigma_\xi=I_n\otimes\xi\quad(\xi\in A).
\]
For the Version~2.3 parameters, $A\cong\F_{q^\ell}$ and
$\dim_{\F_q}A=\ell$.  Multiplication by $\Sigma_\xi$ applies the same
$\ell\times\ell$ matrix $\xi$ independently to every signature block.

The scalar-expanded public base matrices are
$P_1,\ldots,P_{m_1}\in\Mat_{N_0}(\F_q)$.  In the public $q=19$ construction,
each $P_i$ is symmetric.  For $\xi,\eta\in A$, define the bilinear form
\begin{equation}\label{eq:bilinear-form}
 B_i^{\xi,\eta}(x,y)=x^\transpose\Sigma_\xi P_i\Sigma_\eta y.
\end{equation}
Setting $x=y$ makes this a quadratic feature.  The power basis
$1,S,\ldots,S^{\ell-1}$ gives the ordered power-pair coordinates used in the
prior affine-column attack.

\subsection{Imposing one common column variable}\label{sec:affine-columns}

Write the stacked signature columns as
\[
  U=[u_0|\cdots|u_{r-1}],\qquad u_j\in\F_q^{N_0}.
\]
For a matrix block $b\in\{0,\ldots,n-1\}$, our entry convention is
\[
  U_b[a,j]=(u_j)_{b\ell+a},
  \qquad 0\le a<\ell,\quad 0\le j<r.
\]
The attacker chooses multipliers $R_j\in A$ and offsets
$v_j\in\F_q^{N_0}$, with $R_0=1$.  Write
$R=(R_0,\ldots,R_{r-1})$ and
$\mathbf v=(v_0,\ldots,v_{r-1})$, and impose
\begin{equation}\label{eq:relation}
 u_j=\Sigma_{R_j}u+v_j,
 \qquad 0\le j<r.
\end{equation}
Thus all $r$ columns are affine functions of one vector $u$.  This is not an
assumption about honest signatures; it defines a publicly checkable subset in
which the attacker searches for a forgery.

Before the relation is imposed, the feature vector has
\[
  N_{\rm raw}=m_1\ell^2r^2
\]
coordinates: $m_1$ base forms, $\ell^2$ ordered power pairs, and $r^2$ ordered
column pairs.  The output order is public and fixed by the reference
implementation; it is neither secret nor sampled per key.  The implementation
and our equations flatten each $r\times\ell$ output block in opposite orders.
Let $\Pi_{\rm out}$ be the resulting fixed transpose--vectorization
permutation, recorded entry by entry in
\Cref{app:output-permutation}, and put
\[
  t=\Pi_{\rm out}t_{\rm off},\qquad
  \mathcal E=\Pi_{\rm out}\mathcal E_{\rm off}.
\]
This coordinate reordering is valid for square and rectangular parameter sets
and changes neither ranks nor solutions.  In the same convention, the
reference public expansion forms
$\Sigma_{S^a}U$ and then evaluates
$U^\transpose\Sigma_{S^a}P_i\Sigma_{S^b}U$; hence the left--right power
orientation in \eqref{eq:bilinear-form} is the verifier's orientation.

After substituting \eqref{eq:relation}, products of two variable parts give
quadratic terms in $u$, products of one variable part and one offset give
linear terms, and products of two offsets give constants.  For a fixed
relation $R=(R_j)_{j=0}^{r-1}$, let
$\mathcal R_R\in\F_q^{N_{\rm raw}\times m_1\ell^2}$ expand the remaining ordered
quadratic features into the full feature vector, and let
$\mathcal F_{R,\mathbf v}\in\F_q^{N_{\rm raw}\times N_0}$ collect the offset-linear
coefficients.  The complete substituted verifier is
\begin{equation}\label{eq:ordered-system}
 \mathcal E\bigl(\mathcal R_Rz_{\rm ord}(u)+\mathcal F_{R,\mathbf v}u\bigr)
 +c_{R,\mathbf v}=t,
\end{equation}
where $z_{\rm ord}(u)$ contains the $m_1\ell^2$ values
$B_i^{S^a,S^b}(u,u)$ and $c_{R,\mathbf v}$ is the public constant term.  Put
\begin{equation}\label{eq:ER}
 E_R=\mathcal E\mathcal R_R\in\F_q^{M\times m_1\ell^2},
 \qquad \Lambda_{R,\mathbf v}=\mathcal E\mathcal F_{R,\mathbf v}\in\F_q^{M\times N_0}.
\end{equation}

\subsection{Separating quadratic and affine output directions}

By \Cref{thm:symmetric-square}, the ordered feature vector duplicates every
off-diagonal power pair.  Let $z_{\rm sym}(u)$ contain one coordinate for each
unordered pair $0\le a\le b<\ell$, and let $D$ copy each off-diagonal
coordinate into its two ordered positions.  Then
$z_{\rm ord}(u)=Dz_{\rm sym}(u)$.  Define
\begin{equation}\label{eq:QR}
 Q_R=E_RD\in\F_q^{M\times K},
 \qquad K=m_1\binom{\ell+1}{2}.
\end{equation}
The columns of $Q_R$ span all output directions that the residual quadratic
features can reach.

Now choose $W_R$ whose rows form a basis of the left kernel of $Q_R$ and put
\begin{equation}\label{eq:C}
 C_{R,\mathbf v}=W_R\Lambda_{R,\mathbf v}.
\end{equation}
Multiplying by $W_R$ kills every quadratic feature.  What remains is therefore
an ordinary affine system in $u$.  This elementary projection is the key to
retaining all equations of the constrained verifier rather than dropping
output directions that do not lie in the quadratic image.

\begin{center}
\fcolorbox{deepblue}{lightblue}{\parbox{0.92\textwidth}{\small
\textbf{A two-output example.}  Suppose that, after the column substitution,
two verifier outputs over $\F_{19}$ are
\[
  y_1=q(u)+\ell_1(u),\qquad y_2=2q(u)+\ell_2(u),
\]
where $q$ is quadratic and $\ell_1,\ell_2$ are affine.  The combination
$2y_1-y_2=2\ell_1(u)-\ell_2(u)$ cancels the quadratic term.  Solving this
affine equation and substituting it back leaves one quadratic equation; neither
original output was dropped.  The rows of $W_R$ perform the same cancellation
simultaneously for all equations of the constrained verifier.}}
\end{center}

\begin{proposition}[Exact affine reduction]\label{prop:affine-reduction}
Let $\rho=\rank Q_R$, $\lambda=\rank C_{R,\mathbf v}$, and
$\delta=M-\rho-\lambda$.  Then
\[
 \rank[Q_R\ \Lambda_{R,\mathbf v}]=\rho+\lambda,
\]
where $[Q_R\ \Lambda_{R,\mathbf v}]$ denotes horizontal matrix concatenation.  After row
reduction, exactly $\delta$ independent affine consistency conditions remain
on the target.  For every target satisfying those conditions, Gaussian
elimination leaves $\rho$ equations of degree at most two in
$N_0-\lambda$ variables.  Through the affine parametrization, the residual
solution set maps bijectively to the solution set of
\eqref{eq:ordered-system}.

In particular, if $\rho=K$ and $\lambda=M-K$, every target passes the
\emph{affine consistency test}, and the residual system has $K$ equations of
degree at most two on an affine translate of $\ker C_{R,\mathbf v}$.  This conclusion
does not assert that the residual system has a root.
\end{proposition}

\begin{proof}
The row space of $W_R$ is $\ker Q_R^\transpose$.  Since $W_R$ has full row
rank,
\[
 \ker W_R=\bigl(\ker Q_R^\transpose\bigr)^\perp=\operatorname{im}Q_R.
\]
Thus $W_R$ realizes the quotient map
$\F_q^M\to\F_q^M/\operatorname{im}Q_R$.  The image of the columns of
$\Lambda_{R,\mathbf v}$ in this quotient has dimension
$\rank(W_R \Lambda_{R,\mathbf v})=\lambda$, which proves
$\rank[Q_R\ \Lambda_{R,\mathbf v}]=\rho+\lambda$.

Write $b=t-c_{R,\mathbf v}$.  Projecting the substituted verifier by $W_R$ gives
\[
 C_{R,\mathbf v}u=W_R b.
\]
The codomain of $C_{R,\mathbf v}$ has dimension $M-\rho$, while its image has
dimension $\lambda$.  Solvability therefore imposes exactly
$(M-\rho)-\lambda=\delta$ independent affine conditions on the target.
For a target satisfying the affine consistency conditions, parameterize the
solution space of these affine equations by $N_0-\lambda$ free variables.
Choose a linear map $V_R:\F_q^M\to\F_q^\rho$ whose restriction to
$\operatorname{im}Q_R$ is an isomorphism.  Together, $W_R$ and $V_R$ separate
the quotient by $\operatorname{im}Q_R$ from the image itself.  Indeed, if
$T_Rx=(W_Rx,V_Rx)=0$, then $x\in\ker W_R=\operatorname{im}Q_R$, and the
injectivity of $V_R$ on that image gives $x=0$.  The domain and codomain have
the same dimension, so $T_R=(W_R,V_R)$ is an isomorphism.  Substitution
into
\[
 V_R\!\left(Q_Rz_{\rm sym}(u)+\Lambda_{R,\mathbf v}u-b\right)=0
\]
therefore leaves exactly $\rho$ equations of degree at most two, and these
equations together with $C_{R,\mathbf v}u=W_Rb$ are equivalent to all $M$
equations of the substituted verifier.  Thus no row is dropped and the solution set is
preserved bijectively under the affine parametrization.  If $\rho=K$ and $\lambda=M-K$, then $\delta=0$ and
$C_{R,\mathbf v}$ is surjective onto its $M-K$-dimensional codomain; hence
every target passes the affine consistency test.
\end{proof}

\paragraph{Operational meaning.}
The attacker computes $Q_R$, $W_R$, and $C_{R,\mathbf v}$ from the public key before
running an MQ solver.  If the ranks are unfavorable, the relation or offsets
are changed.  When $\rank Q_R=K$ and $\rank C_{R,\mathbf v}=M-K$, every hash target first
determines a nonempty affine space $a(t)+\ker C_{R,\mathbf v}$.  Fix a public
linear right inverse
\[
 J_{R,\mathbf v}:\F_q^{M-K}\longrightarrow\F_q^{N_0},
 \qquad C_{R,\mathbf v}J_{R,\mathbf v}=I,
\]
and take
$a(t)=J_{R,\mathbf v}W_R(t-c_{R,\mathbf v})$.  This base point depends only
on the affine output coordinates, not on the independent residual target.
On that affine space,
$K$ equations of
degree at most two remain.  This is an exact reformulation, not a root-existence
theorem: the residual equations may still have no solution.  Root availability
is analyzed separately in \Cref{sec:slices}.  No verifier equation has been
discarded.  Conversely, every root of the residual system maps through the
chosen column relation to a solution of the complete substituted verifier.

\section{The Symmetry Quotient}\label{sec:symmetry}

The central reduction is the formal version of
\eqref{eq:intro-swap}.  Readers may interpret $\Sym^2(A)$ simply as the vector
space spanned by \emph{unordered} pairs from $A$; the tensor notation records
why the identification is basis-independent.

For a fixed base form $P_i$, define
\[
  q_{i,a,b}(u)=u^\transpose\Sigma_{S^a}P_i\Sigma_{S^b}u.
\]
Transposition gives $q_{i,a,b}=q_{i,b,a}$.  In coordinates, the feature vector
therefore contains two copies of every off-diagonal polynomial.  The theorem
below shows that this is not an accident of the power basis: every alternating
label $\xi\otimes\eta-\eta\otimes\xi$ vanishes.  Write
$\F_q[u]^{\mathrm{hom}}_2$ for the vector space of homogeneous quadratic
polynomials in the coordinates of $u$.

\begin{theorem}[Symmetric-square factorization]\label{thm:symmetric-square}
Let $A=\F_q[S]$ have dimension $d$, and assume $S^\transpose=S$ and
$P_i^\transpose=P_i$.  For each $i$, the map
\[
 A\otimes_{\F_q}A\longrightarrow \F_q[u]^{\mathrm{hom}}_2,
 \qquad
 \xi\otimes\eta\longmapsto
 u^\transpose\Sigma_\xi P_i\Sigma_\eta u
\]
is invariant under $\xi\otimes\eta\mapsto\eta\otimes\xi$.  It therefore
factors through
\[
 \Sym^2_{\F_q}(A)
 =\bigl(A\otimes_{\F_q}A\bigr)
   \big/\Span\{\xi\otimes\eta-\eta\otimes\xi\}.
\]
Consequently, the direct sum over the $m_1$ base forms has rank at most
$m_1\binom{d+1}{2}$.  After any public linear map to $M$ outputs---in
particular, for every relation-dependent matrix $Q_R$---the rank is at most
\begin{equation}\label{eq:symmetry-rank}
 \min\left\{M,m_1\binom{d+1}{2}\right\}.
\end{equation}
If the minimal polynomial of $S$ has degree $\ell$, then $d=\ell$ and, in the
power basis,
\begin{equation}\label{eq:duplication}
 z_{\rm ord}(u)=Dz_{\rm sym}(u).
\end{equation}
\end{theorem}

\begin{proof}
Every element of $A$ is symmetric.  Since the displayed expression is a
scalar,
\[
 u^\transpose\Sigma_\xi P_i\Sigma_\eta u
 =\bigl(u^\transpose\Sigma_\xi P_i\Sigma_\eta u\bigr)^\transpose
 =u^\transpose\Sigma_\eta P_i\Sigma_\xi u.
\]
Thus every alternating label $\xi\otimes\eta-\eta\otimes\xi$ lies in the
kernel.  Taking the direct sum over $i$ gives the first rank bound, and a
linear map to $M$ public outputs gives the additional codomain cap.  When the
powers of $S$ form a basis of $A$, the factorization is exactly the coordinate
identity \eqref{eq:duplication}.
\end{proof}

\paragraph{Why the outer map cannot repair the loss.}
The public outer map receives only the values of these quadratic features.  If
two labeled inputs already define the same polynomial, no later linear
combination can separate them.  The outer map may have full rank on the
$K$-dimensional quotient, but it cannot enlarge that quotient.

For $\ell=4$, each base form loses six of its sixteen ordered coordinates:
$\ell^2=16$ becomes $\binom{\ell+1}{2}=10$.  For $\ell=2$, four ordered
coordinates become three.  These reductions occur for every relation and every
symmetric public key.

The theorem is an upper bound: additional public dependencies could make the
quadratic image even smaller.  The concrete attacks separately check
$\rank Q_R=K$ before solving.  Thus the value $K$ used in the cost tables is
not inferred from symmetry alone; it is the exact rank computed by the public
preflight for the chosen relation.

The factorization itself uses only transposition symmetry and is valid in every
characteristic.  Odd characteristic is needed only for the familiar direct-sum
interpretation
$A\otimes A\cong\Sym^2(A)\oplus\bigwedge^2(A)$.  The submitted $q=16$ system is
outside this theorem because its public base matrices are not symmetrized.

\section{Structured Stable Lifts}\label{sec:structured-lift}

\Cref{prop:affine-reduction} permits arbitrary offsets $v_j$.  Generic offsets
are useful for making the affine equations full rank, but they create linear
terms that mix all coordinates.  That mixing destroys the block grading used
by the selected published $\ell=4$ special-XL model.  We therefore need offsets
that do two jobs at once:
\begin{enumerate}[label=(\roman*),itemsep=2pt]
\item span the missing affine output directions, so that all $M$ equations of
      the substituted verifier are retained; and
\item confine all offset-linear terms to a small public row space, so that a
      large invariant kernel remains available for special XL.\@
\end{enumerate}

The following construction achieves both.  Choose one vector
$w\in\F_q^{N_0}$ and multipliers $\Gamma_0,\ldots,\Gamma_{r-1}\in A$, and set
\begin{equation}\label{eq:structured-offsets}
 v_j=\Sigma_{\Gamma_j}w.
\end{equation}
All offsets are therefore generated from the same $w$ by the public algebra
$A$.  Define the row space
\begin{equation}\label{eq:Uw}
 \mathcal U_w=
 \Span_{\F_q}\left\{
 w^\transpose\Sigma_\xi P_i\Sigma_\eta:
 1\le i\le m_1,\ \xi,\eta\in A
 \right\}
 \subseteq(\F_q^{N_0})^*.
\end{equation}
This space contains every linear form that can arise from crossing a variable
term with one of the structured offsets.  Its a priori bound is $m_1d^2$, not
$m_1\binom{d+1}{2}$: one argument is the fixed offset vector $w$ rather than
the variable $u$ on both sides, so reversing the two algebra labels need not
leave this \emph{linear} form unchanged.

A subspace $H\subseteq\F_q^{N_0}$ is called $A$-linear if
$\Sigma_\xi H\subseteq H$ for every $\xi\in A$.  Since
$A\cong\F_{q^\ell}$, this is the analogue of extension-field linearity.  To
force such a kernel, take the complete $A$-orbit of the affine constraints.
For a basis $1,S,\ldots,S^{d-1}$ of $A$, form
\begin{equation}\label{eq:orbit-matrix}
 \mathcal O_{R,\mathbf v}=
 \begin{bmatrix}
 C_{R,\mathbf v}\\ C_{R,\mathbf v}\Sigma_S\\ \vdots\\ C_{R,\mathbf v}\Sigma_{S^{d-1}}
 \end{bmatrix}.
\end{equation}

\begin{theorem}[Stable-lift kernel]\label{thm:stable-kernel}
Retain the assumptions $S^\transpose=S$ and $P_i^\transpose=P_i$ from
\Cref{thm:symmetric-square}.  For offsets of the form
\eqref{eq:structured-offsets}:
\begin{enumerate}[label=(\roman*)]
\item $\mathcal U_w$ is closed under right multiplication by $A$ and
$\dim_{\F_q}\mathcal U_w\le m_1d^2$;
\item every row of $\mathcal F_{R,\mathbf v}$ and $\mathcal O_{R,\mathbf v}$ lies in
$\mathcal U_w$;
\item if $\rank\mathcal O_{R,\mathbf v}=m_1d^2$, then
\[
 H_{R,\mathbf v}=\ker\mathcal O_{R,\mathbf v}
\]
is $A$-linear, has base-field dimension $N_0-m_1d^2$, and satisfies
$\mathcal F_{R,\mathbf v}H_{R,\mathbf v}=0$.
\end{enumerate}
For the irreducible degree-$\ell$ construction, $H_{R,\mathbf v}$ has $A$-dimension
$n-m_1\ell$.
\end{theorem}

\begin{proof}
See \Cref{app:proof-stable-kernel}.
\end{proof}

\paragraph{Interpretation.}
When the public orbit matrix has the stated rank, its kernel is a large search
direction space closed under $A$.  Moving inside this kernel changes the
quadratic features but changes none of the offset-linear features.  The
offsets can therefore solve the missing affine equations, while the nonlinear
solver sees a clean $A$-linear slice.

\subsection{From an \texorpdfstring{$A$}{A}-linear kernel to four variable blocks}

Special XL needs variables divided into blocks with controlled
multidegrees.  The algebra $A$ supplies those blocks after extending scalars to
a field $\Omega$ in which the minimal polynomial of $S$ splits.  Over
$\Omega$, multiplication by $S$ diagonalizes into $\ell$ eigenspaces.  The
primitive idempotent $e_a$ is simply the projection onto the $a$th eigenspace.
This scalar extension is a solver coordinate change, not a relaxation of the
signature problem: the recovered point must still descend to $\F_q$ and pass
the original verifier.

Let $\mathbf e_a\in\mathbb Z^\ell$ denote the $a$th standard unit vector.  A
quadratic using only block $a$ has multidegree $2\mathbf e_a$; a bilinear term
using blocks $a$ and $b$ has multidegree
$\mathbf e_a+\mathbf e_b$.

\begin{lemma}[Eigenblock coordinates of the quotient]\label{lem:spectral-quotient}
Assume that the minimal polynomial of $S$ is irreducible of degree $\ell$.
Thus $A=\F_q[S]$ is a separable field extension of $\F_q$ of degree $\ell$.
Let $\Omega/\F_q$ be a splitting field of that polynomial, and let
$e_1,\ldots,e_\ell$ be the primitive idempotents of
$A\otimes_{\F_q}\Omega$.  If $L$ is an $A$-linear space of $A$-dimension
$s$, then
\[
 L\otimes_{\F_q}\Omega=L_1\oplus\cdots\oplus L_\ell,
 \qquad L_a=\Sigma_{e_a}(L\otimes_{\F_q}\Omega),
 \qquad \dim_\Omega L_a=s.
\]
For each $i$, the eigenblock features
\[
 x^\transpose\Sigma_{e_a}P_i\Sigma_{e_b}x,
 \qquad 1\le a\le b\le\ell,
\]
are an invertible linear recombination of the symmetric power-pair features.
They have multidegree $2\mathbf e_a$ when $a=b$ and
$\mathbf e_a+\mathbf e_b$ when $a<b$.
\end{lemma}

\begin{proof}
Finite fields are perfect, so the irreducible minimal polynomial of $S$ is
separable.  Over $\Omega$, the power basis and the primitive-idempotent basis
of $A\otimes\Omega\cong\Omega^\ell$ are related by an invertible Vandermonde
matrix.  Passing to the symmetric square gives an invertible change between the
unordered power-pair features and the unordered eigenblock features.  The
idempotent $e_a$ projects $L\otimes\Omega$ onto $L_a$.  Hence a diagonal
feature uses only variables in $L_a$, while an off-diagonal feature is bilinear
in the variables from $L_a$ and $L_b$, giving the stated multidegrees.
\end{proof}

\paragraph{Restricted polynomial rank.}
Write $z_{{\rm sym},1},\ldots,z_{{\rm sym},K}$ for the homogeneous quotient
coordinate polynomials.  For an $\F_q$-linear slice direction $L$, define
\[
 \mathcal Z_L:\F_q^K\longrightarrow\F_q[L]^{\rm hom}_2,
 \qquad
 c\longmapsto\left.(c\cdot z_{\rm sym})\right|_L,
 \qquad
 \rho_{\rm poly}(L)=\rank_{\F_q}\mathcal Z_L.
\]
Thus $\rho_{\rm poly}(L)$ is the dimension of the restricted polynomial
span.  It is distinct from the outer-map rank $\rank_{\F_q}Q_R$, from the
polar rank of an individual scalar quadratic, and from the cross-polar ranks
in \Cref{sec:slices}.  Scalar extension to $\Omega$ preserves
$\rho_{\rm poly}(L)$.

\begin{corollary}[Stable affine slices]\label{cor:stable-slices}
Assume that $q$ is odd, that the minimal polynomial of $S$ is irreducible of
degree $\ell$, and that the public base-field rank checks
\begin{equation}\label{eq:stable-preflight}
 \rank_{\F_q} Q_R=K,\qquad
 \rank_{\F_q} C_{R,\mathbf v}=M-K,\qquad
 \rank_{\F_q}\mathcal O_{R,\mathbf v}=m_1\ell^2
\end{equation}
pass, and let $L\subseteq H_{R,\mathbf v}$ be an $A$-linear subspace of
$A$-dimension $s$ (equivalently, base-field dimension $h=\ell s$).  Assume
also that the restricted polynomial-rank check
$\rho_{\rm poly}(L)=K$ passes.  For every
target, the affine equations in \Cref{prop:affine-reduction} have a
solution space $a(t)+\ker C_{R,\mathbf v}$ for a publicly computed base point
$a(t)$.  For every $y\in\ker C_{R,\mathbf v}$, the affine
slice $a(t)+y+L$ is contained in that space, and the full offset-linear feature
$\mathcal F_{R,\mathbf v}u$ is constant on the slice.

After extension to a splitting field, use the eigenblock coordinates of
\Cref{lem:spectral-quotient}.  Translation by $a(t)+y$ creates only linear and
constant terms in the same one or two blocks as the corresponding quadratic
term.  Introducing one homogenizing variable per block---an auxiliary variable that
is set to one after solving---therefore turns each complete affine equation
into a homogeneous equation of multidegree
$2\mathbf e_a$ or $\mathbf e_a+\mathbf e_b$.  The disjoint monomial supports
of these unordered multidegrees force $\Omega$-rank $m_1$ in every component:
there are
$\binom{\ell+1}{2}$ components, each has rank at most $m_1$, and their ranks
sum to $\rho_{\rm poly}(L)=K=m_1\binom{\ell+1}{2}$.
\end{corollary}

\begin{proof}
See \Cref{app:proof-stable-slices}.
\end{proof}

\paragraph{What the preflight certifies.}
The first two ranks in \eqref{eq:stable-preflight} say that the quotient has
its full expected rank and that the offsets supply every remaining affine
output direction.  The orbit rank says that the kernel has the predicted
$A$-linear dimension and kills all offset-linear motion.  The restricted
polynomial-rank check $\rho_{\rm poly}(L)=K$ then certifies that no
multidegree component was lost on the chosen slice.  All of these tests are
public and occur before the expensive Macaulay computation.  A failing
relation, offset family, or slice is rejected and the preflight is repeated
with another public choice; the evidence in \Cref{sec:evidence} records the
observed number of choices, but the theorem does not guarantee that this
search succeeds for every valid key.

\section{Verifier Rejection Compatibility}\label{sec:rejection}

The public verifier does more than evaluate the quadratic map.  It also rejects
signatures containing too many symmetric blocks.  A root of the reduced MQ
system is useful only if it passes this block-symmetry rejection test.  We therefore build
rejection compatibility into the reduction rather than treating it as an
unmodeled success probability.

Version~2.3 rejects a square odd-characteristic signature if one of its blocks
is symmetric when $\ell>2$, and rejects an $\ell=2$ signature if more than
$\lfloor n/4\rfloor$ blocks are symmetric~\cite{SNOVA23}.  Rectangular blocks
are not subject to this rule.

\paragraph{The structured $\ell=4$ rows.}
For each square block $b$, substituting \eqref{eq:relation} into the equalities
$U_b[a,c]=U_b[c,a]$ gives an affine system
\begin{equation}\label{eq:rejection-affine}
 G_bu=-h_b.
\end{equation}
Here $G_b$ is the public coefficient matrix of the variable part of those
equalities, and $h_b$ is their public offset vector.
If $\rank[G_b\mid-h_b]>\rank G_b$, this system is inconsistent, so no point in
the entire residual search space can make block $b$ symmetric.  Requiring this
inconsistency for every square block is an exact public certificate that every
root produced by the structured lift passes that check.

\paragraph{A deterministic $\ell=2$ construction.}
All three parameter levels use offsets that force a fixed nonzero skew in
every block.

\begin{proposition}[Fixed-skew $\ell=2$ lift]\label{prop:fixed-skew}
For the Version~2.3 $\ell=2$ matrix
\[
 S=\begin{pmatrix}1&2\\2&15\end{pmatrix},
 \qquad R_0=I,
 \qquad R_1=9I+10S=\begin{pmatrix}0&1\\1&7\end{pmatrix},
\]
choose offsets satisfying
\begin{equation}\label{eq:fixed-skew-offset}
 (v_1)_{b,0}-(v_0)_{b,1}=1
\end{equation}
for every block $b$.  Then every block of every signature satisfying the
column relation has
\begin{equation}\label{eq:fixed-skew-identity}
 U_b[0,1]-U_b[1,0]=1.
\end{equation}
Consequently every solution passes the $\ell=2$ block-symmetry rejection
test.  If the public preflight also gives
$\rank Q_R=K$ and $\rank C_{R,\mathbf v}=M-K$, every target passes the
affine consistency test and leaves $K$ equations of degree at most two in
$N_0-M+K$ variables.  Every root of that residual system reconstructs to a
root of the substituted verifier that passes the block-symmetry rejection
test; after encoding, it must still pass the complete unmodified verifier.
\end{proposition}

\begin{proof}
The first row of $R_1$ is the second row of $R_0=I$.  Hence the variable part
of $U_b[0,1]-U_b[1,0]$ vanishes.  The remaining offset difference is one by
\eqref{eq:fixed-skew-offset}, proving \eqref{eq:fixed-skew-identity}.  No block
is symmetric.  The final statement is the full-rank case
$\rho=K$, $\lambda=M-K$ of \Cref{prop:affine-reduction}.
\end{proof}

Thus every fixed-skew MQ root passes the block-symmetry rejection test
deterministically.  Serialization, canonical encoding,
message--salt freshness, and every other public check are still applied by the
final unmodified verifier.  All three $\ell=2$ headline rows use this
fixed-skew lift.

A lower-per-solve-cost zero-offset alternative has a nontrivial acceptance
problem: its roots may contain too many symmetric blocks.  The exact
accepted-root bounds are collected in \Cref{app:zero-offset}.  They prove
existence and population-fraction statements under polar-rank hypotheses, but
do not determine which root a deterministic or nonuniform MQ solver returns.
Accordingly, this alternative is not used in the headline table; its numerical
sensitivity analysis also appears in \Cref{app:zero-offset}.

\section{Root Existence, Uniqueness, and Retry Bounds}\label{sec:slices}

After the exact affine reduction, the structured attack repeatedly chooses a
hash target and an affine coset of the selected slice direction space.  We call
one such pair a \emph{target--slice trial}, equivalently a target--coset pair.
The primer in \Cref{sec:mq-primer} distinguishes root existence, uniqueness,
and solver readout.  This section treats the first two; the third remains the
separate special-XL solving assumption in \Cref{sec:special-xl}.

Concretely, after the affine target coordinates determine a base point
$a(t)$, the attack samples $y+L$ uniformly from
$\ker C_{R,\mathbf v}/L$, independently of the retained residual target and
of every other fresh trial.  The target-dependent translation $a(t)+y$ does
not change any polar map.  This is the sampling rule used when the abstract
theorems below take a uniform coset of $V/L$.

Let $V$ be the full residual direction space, let $L\subseteq V$ be the chosen
slice direction space, and write $h=\dim L$.  For a map of degree at most two
$g:V\to\F_q^K$, a random-map heuristic predicts
\[
  \mu=q^{h-K}
\]
roots in a random affine coset of $L$ above a random target.  When $h=K-\tau$ with $\tau\ge0$, this mean is $q^{-\tau}$, suggesting
about $q^\tau$ trials.  The obstruction detected by a second-moment analysis is excess
collisions: distinct points in the same slice mapping to the same target.
Polar derivatives count those collisions exactly, so the argument below is a
theorem rather than a random-system heuristic.  The cleanest case is full
cross-polar rank: then distinct points behave almost as independently as the
mean-root calculation suggests.  For the positive slice deficits used here, a
unique-root trial costs essentially $q^{K-h}=q^\tau$ attempts.  The theorem first states the exact formulas and
then quantifies what changes when some output directions lose rank.

Write the vector-valued polar map as
\[
 \beta(x,y)=g(x+y)-g(x)-g(y)+g(0).
\]
For $b\in\F_q^K\setminus\{0\}$, define the cross-polar map
\begin{equation}\label{eq:cross-polar-map}
 \Phi_{b,L}:V\longrightarrow L^*,
 \qquad \Phi_{b,L}(x)(y)=b\cdot\beta(x,y),
\end{equation}
and let $e_L(b)=\rank\Phi_{b,L}$.  Put
\begin{equation}\label{eq:theta-def}
 e_* = \min_{b\ne0}e_L(b),
 \qquad
 \Theta_L(g)=\sum_{b\ne0}q^{-e_L(b)}.
\end{equation}
For $y\in L\setminus\{0\}$, also define the directional derivative
\begin{equation}\label{eq:directional-map}
 T_y:V\longrightarrow\F_q^K,
 \qquad T_y(x)=\beta(x,y),
\end{equation}
write $r(y)=\rank T_y$, and let
$\kappa(y)$ be one when $-(g(y)-g(0))\in\operatorname{im}T_y$ and zero
otherwise.  The indicator $\kappa(y)$ records whether two points separated by
$y$ can collide at all; the rank $r(y)$ then determines how many such pairs
exist.

\begin{theorem}[Root collisions and the rank-spectrum bound]
\label{thm:coset-success}
Choose an affine coset $C$ uniformly from $V/L$ and choose
$z\in\F_q^K$ uniformly.  Let
\[
 Z=|\{x\in C:g(x)=z\}|,
 \qquad \mu=q^{h-K},
\]
and define
\begin{equation}\label{eq:eta-def}
 \eta_0=\sum_{y\in L\setminus\{0\}}\kappa(y)q^{-r(y)},
 \qquad
 \bar\eta=\sum_{y\in L\setminus\{0\}}q^{-r(y)}.
\end{equation}
Then
\begin{align}
 \mathbb E Z&=\mu,\label{eq:EZ}\\
 \mathbb E[Z(Z-1)]&=\mu\eta_0,\label{eq:factorial-moment}\\
 \eta_0&\le\bar\eta
 =\mu\bigl(1+\Theta_L(g)\bigr)-1.\label{eq:eta-spectrum}
\end{align}
Consequently,
\begin{align}
 \operatorname{Var}(Z)&\le\mu^2\Theta_L(g),\label{eq:theta-var}\\
 \Pr[Z>0]&\ge\frac{\mu}{1+\eta_0}
 \ge\frac{1}{1+\Theta_L(g)}.\label{eq:coset-success}
\end{align}
If $h=K-\tau$ and $e_*\ge h-\Delta$, then
\begin{equation}\label{eq:defect-retry}
 \Theta_L(g)<q^{\tau+\Delta},
 \qquad
 \mathbb E[\textnormal{independent trials to a root-containing
 target--coset pair}]
 <1+q^{\tau+\Delta}.
\end{equation}
This last bound controls the search for a root-containing target--coset pair,
not uniqueness of the root.
Here $\Delta$ is the maximum loss from full cross-polar rank $h$.
\end{theorem}

\begin{proof}
See \Cref{app:proof-coset-success}.
\end{proof}

\paragraph{How to read the theorem.}
The first moment $\mathbb EZ=\mu$ is exactly the random-map value.  The second
factorial moment $\mathbb E[Z(Z-1)]$ counts ordered pairs of distinct roots and
therefore measures collisions.  The compatible collision mass $\eta_0$ is
exact; $\bar\eta$ drops the compatibility test and is a convenient upper
bound.  The sum $\Theta_L(g)$ describes the entire cross-polar rank spectrum,
not merely its minimum.  A few isolated rank defects may have little effect,
whereas low rank in many output directions can substantially increase retries.

\begin{center}
\begin{tabularx}{0.92\textwidth}{@{}>{\bfseries}p{0.27\textwidth}X@{}}
$\mu=q^{h-K}$ & expected number of roots in a random target--slice pair\\
$\eta_0$ & exact expected number of compatible collisions per planted root\\
$\bar\eta$ & rank-only upper bound on that collision mass\\
full cross-polar rank & for the selected positive deficits, unique-root trial
count essentially $q^{K-h}=q^\tau$
\end{tabularx}
\end{center}

\begin{corollary}[Unique-root trials and the planted distribution]
\label{cor:singleton-coset}
In the setting of \Cref{thm:coset-success},
\begin{align}
 \Pr[Z=1]&\ge\mu(1-\eta_0)
             \ge\mu(1-\bar\eta),\label{eq:singleton-coset}\\
 \Pr[Z>0]&\ge\mu\left(1-\frac{\bar\eta}{2}\right),
 \qquad (\bar\eta<2),\label{eq:bonferroni-success}\\
 \Pr[Z\ge2\mid Z>0]&\le\frac{\bar\eta}{2-\bar\eta},
 \qquad (\bar\eta<2).\label{eq:conditional-multiple}
\end{align}
In particular, if $\bar\eta<1$, the expected number of independent trials
until $Z=1$ is at most
\begin{equation}\label{eq:unique-retry-general}
 \frac{1}{\mu(1-\bar\eta)}.
\end{equation}

Let $\mathsf U$ be the uniform target--coset pair conditioned on $Z>0$, and
let $\mathsf P$ be the planted distribution obtained by choosing a uniform
coset, a uniform point in it, and the point's image as target.  Conditioned on
$Z=1$, both laws equal the uniform target--coset law conditioned on $Z=1$.  If
$\bar\eta<1$, then
\begin{equation}\label{eq:planted-tv}
 \Pr_{\mathsf P}[Z\ge2]\le\eta_0\le\bar\eta,
 \qquad
 d_{\rm TV}(\mathsf U,\mathsf P)\le\bar\eta.
\end{equation}
\end{corollary}

\begin{proof}
See \Cref{app:proof-singleton-coset}.
\end{proof}

Total variation distance is the largest discrepancy that the two distributions
can assign to any event.  Thus, when $\bar\eta$ is tiny, an experiment that
plants a known root samples essentially the same target--coset distribution as
a uniformly chosen trial conditioned on being solvable.  This justifies using
planted instances to probe solver behavior; it does not remove the separate
finite-size-to-production extrapolation.

\begin{corollary}[Full cross-polar rank]\label{cor:full-cross-rank}
The following are equivalent:
\begin{enumerate}[label=(\roman*)]
\item $e_L(b)=h$ for every $b\ne0$;
\item $T_y$ is surjective for every $y\in L\setminus\{0\}$.
\end{enumerate}
Under these conditions,
\[
 \eta_0=\bar\eta=\mu-q^{-K}.
\]
If $h=K-\tau$ with $\tau\ge0$, the expected number of trials until a unique
base-field root is at most
\begin{equation}\label{eq:unique-retry}
 \mathsf{Retry}_{K,\tau}\coloneqq
 \frac{q^\tau}{1-q^{-\tau}+q^{-K}}.
\end{equation}
\end{corollary}

\begin{proof}
If some $T_y$ is not surjective, a nonzero $b$ annihilates its image, so the
transpose of $\Phi_{b,L}$ has the nonzero kernel vector $y$ and $e_L(b)<h$.  The
converse is the same argument in reverse.  Under either equivalent condition,
$r(y)=K$ and $\kappa(y)=1$ for every nonzero $y$, giving
$\eta_0=\bar\eta=(q^h-1)q^{-K}=\mu-q^{-K}$.  Substitute this value into
\eqref{eq:unique-retry-general}.
\end{proof}

For the positive deficits used in the concrete attacks, the expected
unique-root trial count under full cross-polar rank is essentially $q^\tau$:
the denominator in \eqref{eq:unique-retry} differs from one only by the tiny
terms $q^{-\tau}$ and $q^{-K}$.  In the running Level-I example, $h=40$, $K=50$,
and $\tau=10$, so the expected number of independent target--slice trials is
essentially $19^{10}=2^{42.48}$.  Conditional on a root-containing trial, the
probability of more than one base-field root is below $2^{-43.47}$.  This is
the precise justification for the retry factor used in the structured cost
model.

Equation~\eqref{eq:eta-spectrum} distinguishes sparse rank defects from a
pervasive loss.  If at most $T$ nonzero coefficient vectors $b$ are defective
and each has defect $h-e_L(b)\le\Delta$, then
\begin{equation}\label{eq:defect-spectrum}
 \bar\eta\le
 \mu-q^{-K}+Tq^{-K}(q^\Delta-1).
\end{equation}
If defects are counted projectively, $T$ in this formula must include all
$q-1$ nonzero scalar multiples of each projective direction.  Under full
cross-polar rank and $q=19$, \eqref{eq:conditional-multiple} is below
$2^{-43.47}$ for $\tau=10$ and below $2^{-60.47}$ for $\tau=14$; the
corresponding total-variation bounds are below $2^{-42.47}$ and $2^{-59.47}$.

\begin{corollary}[Full-space target availability]\label{cor:full-space}
Let $q$ be odd and let $g:\F_q^N\to\F_q^K$ have degree at most two.  Let
$r_*$ be the
minimum polar rank of a nonzero output combination, and put
$\alpha=1-q^{-K}$.  For a uniform target,
\begin{equation}\label{eq:full-space-success}
 \Pr[g^{-1}(z)\ne\varnothing]
 \ge \frac{1}{1+\alpha q^{K-r_*}}.
\end{equation}
Hence, for independent repetitions, the expected number of target trials is at most
$1+\alpha q^{K-r_*}$.  If $r_*\ge2K$, every target has a root.
\end{corollary}

\begin{proof}
See \Cref{app:proof-full-space}.
\end{proof}

For a structured $\ell=4$ lift, an invertible output transformation separates
the $M-K$ affine coordinates from the remaining $K$ coordinates of a uniform
target.  Conditioning on the affine coordinates fixes a base point; the other
coordinates remain uniform, and translation does not change the polar maps.
We take $V=\ker C_{R,\mathbf v}$ and let $L$ be the selected stable slice, of dimension
$h=\ell s=K-\tau$.  Full cross-polar rank gives the unique-root retry bound
$\mathsf{Retry}_{K,\tau}$ in \eqref{eq:unique-retry}.  A worst-case
rank-defect bound gives the budget for reaching a root-containing target--coset
pair in \eqref{eq:defect-retry}; unique recovery instead
depends on the aggregate quantity $\bar\eta$ and the Macaulay kernel model.

For the fixed-skew $\ell=2$ lift, take $L=V$ to be the complete residual space.
By \Cref{prop:fixed-skew}, every root passes the block-symmetry rejection test.  The
expected target-search
factor is at most $1+\alpha q^{K-r_*}$, and $r_*\ge2K$ makes every target
solvable.  The concrete rank thresholds, retry penalties, and evidence are
collected in \Cref{sec:l2-costs,tab:fixed-skew-evidence}; all are global
conditions on nonzero output directions, not consequences of sampled ranks.

\subsection{Uniform targets in the random-oracle experiment}
\label{sec:target-sampling}

The probability theorems above assume a uniform target in $\F_q^K$.
This subsection states precisely the idealization under which the official
hash-to-field conversion supplies that distribution.  We model the
domain-separated SHAKE calls in the nested target computation as random
oracles.  A \emph{fresh trial} is one whose actual outer-oracle input has not
appeared in any earlier trial.  The forger skips any duplicate.  Conditional
on freshness, the binary chunks returned on distinct trials are independent
and uniform.  The statements below are exact within this random-oracle
experiment; they are not standard-model statements about deterministic SHAKE.

Concretely, the official $q=19$ conversion reads at most fifteen base-19 digits
from each 64-bit chunk.  If a uniform $b$-bit chunk $X$ is to supply $a$
digits, the attacker accepts that message--salt candidate only if
\begin{equation}\label{eq:uniform-chunk}
 X<\left\lfloor\frac{2^b}{q^a}\right\rfloor q^a.
\end{equation}
Conditioned on this event, $X\bmod q^a$ is uniform, so the extracted
digits are jointly uniform.  For independent chunks, conditioning on the
product event that \eqref{eq:uniform-chunk} holds for every chunk therefore
makes the retained official target exactly uniform in $\F_{19}^M$.  The fixed
invertible output-coordinate change in \Cref{sec:affine-columns} preserves
uniformity.  Moreover, the proof of \Cref{prop:affine-reduction} shows that
\[
 T_R=(W_R,V_R):\F_q^M\longrightarrow\F_q^{M-K}\times\F_q^K
\]
is an isomorphism when $\rank Q_R=K$.  Hence
$T_R(t-c_{R,\mathbf v})$ is uniform, and conditioning on its affine
coordinates leaves its residual $K$ coordinates uniform and independent.
Condition on a value of the affine block
$W_R(t-c_{R,\mathbf v})$.  This fixes the base point $a(t)$ chosen in
\Cref{sec:prelim} and hence fixes the translated residual map $g$.  Given that
block, the residual target is still uniform, while the affine coset is sampled
uniformly and independently from $\ker C_{R,\mathbf v}/L$.  The conditional
law is therefore exactly the law in \Cref{thm:coset-success}.  Translation
does not change the polar maps, so the stated rank hypotheses and bounds are
uniform in the conditioned affine block; averaging over that block preserves
them.  Successive fresh trials are independent within the random-oracle
experiment.  This is
target \emph{preselection}; it is distinct
from testing whether a solver-produced root passes the block-symmetry
rejection test.

The forger performs this rejection sampling over message--salt candidates;
the verifier continues to apply the unmodified official conversion.  Across
all Version~2.3 target lengths, the hash-chunk acceptance probability is at
least $0.152462$, or fewer than 6.6 candidates on average.  Distinct
message--salt pairs need not produce distinct outer inputs because the message
is first compressed to a 512-bit digest.  Freshness can nevertheless be
enforced with a simple schedule: for a fixed chosen-message prefix, enumerate
salts without repetition; before beginning a new prefix, compute its digest
and discard the prefix if that digest has appeared before.  The resulting
$(\text{digest},\text{salt})$ pairs are pairwise distinct.  This makes the
modeled independence exact even if an inner-digest collision occurs, rather
than merely treating such a collision as negligible.

The larger salt-space and prefix budget needed by the unselected zero-offset
Level-V alternative is separated from this general sampling argument and
accounted for in \Cref{app:zero-offset}.

\section{Conditional Arithmetic Cost Estimates}\label{sec:complexity}

The algebraic reduction determines the residual MQ systems exactly.  Turning
those systems into numerical security estimates requires a separate cost
model.  We report that model in layers so that a reader can see which numbers
are structural, which come from an MQ solver estimate, and which are merely a
change of arithmetic units.

\begin{table}[H]
\centering
\caption{How to interpret the cost accounting.}
\label{tab:cost-layers}
\small
\renewcommand{\arraystretch}{1.12}
\begin{tabularx}{0.96\textwidth}{@{}>{\bfseries}p{0.23\textwidth}X@{}}
\toprule
Layer & What it means\\
\midrule
Residual system & Exact numbers of equations, variables, blocks, and public
rank conditions after the reduction.\\
Solver model & Published generic-MQ or special-XL formula for the dominant
finite-field work at those dimensions.\\
AXN conversion & Replacement of the solver formula's coarse field-operation
factor by an explicit two-input AND/XOR/XNOR circuit for one signed
multiply-accumulate.\\
Not included & A complete implementation of sparse matrix construction,
pivoting and control, memory addressing and traffic, communication, or an
optimized end-to-end attack netlist.\\
\bottomrule
\end{tabularx}
\end{table}

A cost exponent $c$ means modeled work proportional to $2^c$ in the stated
unit.  A difference of ten exponent bits is a factor of about $2^{10}$, not ten
percent.  Because the solver formulas count dominant arithmetic rather than a
complete machine, we call the difference from a NIST reference an
\emph{nominal arithmetic gap}.

\subsection{Cost units and circuit conversion}\label{sec:units}

The cited attacks first count finite-field operations and then apply coarse
bit-cost factors~\cite{Beullens2024SNOVA,Cabarcas2025SNOVA}.  NIST lists
$2^{143}$, $2^{207}$, and $2^{272}$ as nominal classical-gate references, while
also requiring category comparisons across every metric it considers relevant
to practical security~\cite{NISTPQCCriteria}.  We retain each attack's
field-operation count and replace only its dominant field factor by explicit
circuits.  The resulting figures sharpen one metric; they do not by themselves
establish a category break.  Our primary AXN basis assigns unit cost to two-input
AND, XOR, and XNOR gates, with free constants, wires, fan-out, and structural
sharing.  This is the gate family used by the concrete AES-128 circuit on
NIST's circuit list~\cite{NISTCircuitComplexity}.  A two-input-NAND conversion
provides a separate basis-sensitivity check.

For generic MQ, the published factor is
\begin{equation}\label{eq:published-f19-factor}
 b_1=(\log_2 19)^2+\log_2 19=22.29.
\end{equation}
The validated circuits described in
\Cref{app:field-circuits} give the conservative signed
multiply-accumulate charges
\begin{equation}\label{eq:axn-f19-factor}
 g_1^{\mathrm{AXN}}=441,
 \qquad g_1^{\mathrm{NAND}}=835.
\end{equation}
For the $\F_{19^4}$ arithmetic used by special XL, the corresponding charges
are
\begin{equation}\label{eq:axn-f194-factor}
 g_4^{\mathrm{AXN}}=6{,}081,
 \qquad g_4^{\mathrm{NAND}}=11{,}102.
\end{equation}
These values replace only the local arithmetic primitive charged by the
published estimators.  They do not upper-bound a complete solver.

\paragraph{Same-basis calibration.}
Across the nine rows, the AXN arithmetic schedules are
\AXNExpectedAESGapRange{} bits below the reported expected AES key-search
calibrations, \AXNFullAESGapRange{} bits below full-key-space AES enumeration,
and \AXNNominalGapRange{} bits below the nominal NIST references.  The NAND
sensitivity ranges are \LoneNANDRange, \LthreeNANDRange, and
\LfiveNANDRange.  The circuit constructions, validation checks, and AES table
are deferred to \Cref{app:field-circuits,app:aes-calibration}.  These are
resource-asymmetric calibrations, not end-to-end attack comparisons.

\paragraph{Scope of the conversion.}
Target generation is treated separately in \Cref{sec:target-sampling};
public preprocessing is not priced.  The dominant unresolved inputs are the
global production rank spectra and the accuracy of the production-scale
solving models; the implementation costs excluded in
\Cref{tab:cost-layers} remain outside the conversion.

\subsection{Structured special XL for \texorpdfstring{$\ell=4$}{l=4}}\label{sec:special-xl}

The exact reduction supplies an $A$-linear slice of $A$-dimension $s$,
hence base-field dimension $4s$ when $\ell=4$.  After extension to the
splitting field, this becomes four blocks of $s$ variables.  For each block
pair $1\le a\le b\le4$, there are $m_1$ independent equations
of degree two \emph{when the public restricted-polynomial-rank preflight
$\rho_{\rm poly}(L)=K$ passes}: diagonal equations use one block twice, and off-diagonal
equations use two blocks once each.  Special XL records a separate degree for
each block and selects a multidegree
at which the associated Macaulay matrix is predicted to have the needed rank.
A reader need not manipulate the series below: here it is only a compact device
for choosing the matrix degree and counting the resulting monomial columns.

After adding one homogenizing variable to each block, the Hilbert series used
by the published special-XL model is
\begin{equation}\label{eq:hilbert-series}
 H_{m_1,s}(\mathbf t)=
 \frac{\prod_{a=1}^{4}(1-t_a^2)^{m_1}
       \prod_{1\le a<b\le4}(1-t_at_b)^{m_1}}
      {\prod_{a=1}^{4}(1-t_a)^{s+1}}.
\end{equation}
For $\mathbf d=(d_1,\ldots,d_4)$, put
\begin{equation}\label{eq:macaulay-size}
 M_s(\mathbf d)=\prod_{a=1}^{4}\binom{s+d_a}{d_a}.
\end{equation}
The search chooses a multidegree whose signed Hilbert coefficient is
negative.  This sign change is a rank-prediction trigger in the cited
special-XL model: a random Macaulay matrix is predicted to have full column
rank, while a system with a planted solution is predicted to have a
one-dimensional right kernel
\cite{Cabarcas2025SNOVA}.  The modeled linear-algebra work for one target--slice trial is
$3s^2\bigl(M_s(\mathbf d)\bigr)^2$ extension-field multiplications.  This is
the quadratic linear-algebra schedule adopted in the cited special-XL
estimator---specifically its $3\max_a(s_a)^2M(\mathbf d)^2$,
matrix-multiplication-exponent-$2$ convention---not a claim that ordinary
dense Gaussian elimination is quadratic~\cite{Cabarcas2025SNOVA}.  The factor
$3s^2$ is the estimator's stated prefactor.

Cabarcas et al.\ handle an underdetermined system by specializing $k$
coordinates, where $4\mid k$, and multiplying the work by $q^k$.  Our attack
uses a different retry loop: it fixes a $4s$-dimensional slice and samples the
target and affine coset.  Under full cross-polar rank,
\Cref{cor:full-cross-rank} bounds the expected trials needed for a unique
base-field root by $\mathsf{Retry}_{K,\tau}$.  The modeled work is therefore
\begin{equation}\label{eq:special-xl-cost}
 \mathsf{Retry}_{K,\tau}\,3s^2M_s(\mathbf d)^2,
 \qquad \tau=K-4s,
\end{equation}
extension-field multiplications.  Replacing the approximation $1+q^\tau$ by
the exact $\mathsf{Retry}_{K,\tau}$ bound changes each displayed logarithm by
less than $10^{-25}$ bits.
If every cross-polar map has rank defect at most $\Delta$, then
\eqref{eq:defect-retry} bounds the expected target--coset trials---and hence
the associated Macaulay linear-algebra budget under the same per-trial
model---by $1+q^{\tau+\Delta}$.  It does not by itself imply unique
recovery.  That step depends on $\bar\eta$ and
on the Macaulay kernel behavior described below.

The computation takes place over $\F_{19^4}$.  The published convention assigns
each extension-field multiplication
\begin{equation}\label{eq:bitop-conversion}
 b_4=2\bigl(\log_2 19^4\bigr)^2+\log_2 19^4=594.43
\end{equation}
bit operations.  The AXN arithmetic conversion replaces $b_4$ by the conservative
signed multiply-accumulate bound $g_4^{\mathrm{AXN}}$ from
\eqref{eq:axn-f194-factor}.

\begin{lemma}[Recovery from a one-dimensional Macaulay kernel]\label{lem:macaulay-recovery}
Fix a target $t$, a selected affine slice $a(t)+y+L$, and an
$\F_q$-linear parameterization $\iota:\F_q^h\to L$.  Let $\lambda_*\in\F_q^h$ be slice coordinates for a
root
\[
  u_*=a(t)+y+\iota(\lambda_*),
\]
and let $x_*\in\Omega^h$ be the image of $\lambda_*$ under the eigenblock
change of variables.  Let $\mathcal M$ be the corresponding special-XL
Macaulay matrix over the splitting field $\Omega$, dehomogenized by setting
every block homogenizer to one.  Suppose its columns include the constant and
degree-one monomials.  If $\ker_\Omega\mathcal M$ is one-dimensional, row
reduction recovers $x_*$, $\lambda_*$, the common column $u_*$, and hence the
complete signature candidate.
\end{lemma}

\begin{proof}
The dehomogenized monomial-evaluation vector $\nu(x_*)$ lies in
$\ker_\Omega\mathcal M$ and has constant coordinate one.  Every nonzero
kernel vector is therefore a scalar multiple of $\nu(x_*)$.  Normalizing its
constant coordinate to one recovers $x_*$ from the degree-one coordinates.  The
inverse eigenblock transformation recovers $\lambda_*\in\F_q^h$; the affine
embedding gives $u_*=a(t)+y+\iota(\lambda_*)$; and the column relation
\eqref{eq:relation} reconstructs every signature column.
\end{proof}

\paragraph{Las Vegas wrapper.}
For each target--slice trial, compute the Macaulay kernel and continue unless
it is one-dimensional with nonzero constant coordinate.  Normalize that
coordinate, recover the degree-one eigenblock coordinates, and invert the
eigenblock transformation.  Continue again unless the point lies in $\F_q$
and the public verifier accepts it.  The procedure therefore never returns an
invalid signature; the solving model predicts only how often a trial reaches
the recoverable case with a one-dimensional right kernel.

\paragraph{Recovery probability and sensitivity.}
The root theorem counts $\F_q$-rational roots, so a unique-root trial contains
the point required by \Cref{lem:macaulay-recovery}.  The remaining assumption
that the production Macaulay right kernel is
one-dimensional concerns extra algebraic kernel vectors, including those
caused by extension-field roots or by the Macaulay rows failing to span all
expected relations; it is not an uncharged descent retry.  The small-profile
tests in \Cref{sec:evidence} support, but do not prove, this production model,
and every recovered point is checked by the original verifier.

When $\bar\eta<1$, this remaining dependence can be parameterized without
conflating solvability with uniqueness.  Let $\mathsf P_1$ be the planted distribution conditioned
on $Z=1$.  By \Cref{cor:singleton-coset}, this is exactly the same law as a
uniform target--slice trial conditioned on $Z=1$.  Let $\mathcal R$ be the
event that the chosen Macaulay matrix has the one-dimensional,
nonzero-constant-coordinate kernel required by
\Cref{lem:macaulay-recovery}, and put
\[
 p_{\rm ker}^{(1)}=\Pr_{\mathsf P_1}[\mathcal R].
\]
If $p_{\rm ker}^{(1)}>0$, then by \eqref{eq:singleton-coset}, a fresh trial
reaches the recoverable event with probability at least
\[
 \mu(1-\bar\eta)\,p_{\rm ker}^{(1)}.
\]
For independent repetitions, its expected retry count is therefore at most
\[
 \frac{1}{\mu(1-\bar\eta)p_{\rm ker}^{(1)}}.
\]
The displayed special-XL schedules use the conditional production model
$p_{\rm ker}^{(1)}=1$; a justified lower bound below one adds
$-\log_2p_{\rm ker}^{(1)}$ bits.  Thus lower bounds $0.9$, $0.5$, and $0.1$
would add about $0.152$, $1$, and $3.322$ bits, respectively.  A point
estimate from the small campaign is not such a lower bound.  If the kernel has
dimension greater than one, the readout lemma
does not apply: one must instead price additional root enumeration, a
change-of-ordering algorithm, or a border-basis computation
\cite{FaugereEtAl1993FGLM,KehreinKreuzer2005}.
There is no justified generic ``multiply by the nullity'' correction for
nullity greater than one because a kernel basis need not enumerate evaluation
vectors.  Likewise, if the Hilbert-series sign change occurs earlier or later
than modeled, one must replace $\mathbf d$ by the observed stopping
multidegree and rescale the per-trial term by
$\bigl(M_s(\mathbf d_{\rm obs})/M_s(\mathbf d)\bigr)^2$ before applying the
retry factors.

The displayed schedules use global full cross-polar rank and the conditional
modeled value $p_{\rm ker}^{(1)}=1$.  A rank-defect bound alone controls
retries to a
root-containing target--slice trial through \eqref{eq:defect-retry}; it does
not imply unique recovery or a one-dimensional Macaulay kernel.

\begin{table}[tbp]
\centering
\caption{Cost-driving special-XL parameters for the six $\ell=4$ rows,
conditional on global full cross-polar rank and the modeled value
$p_{\rm ker}^{(1)}=1$.  ``Model'' uses $b_4$; ``AXN arithmetic'' replaces it by the
author-reported signed multiply-accumulate circuit
$g_4^{\mathrm{AXN}}$.}
\label{tab:xl-optima}
\scriptsize
\setlength{\tabcolsep}{4pt}
\begin{tabular}{@{}c c c c c r r@{}}
\toprule
$m_1$ & available $\dim_A H$ & $K$ & $(s,\tau)$ & $\mathbf d$ & model & AXN arithmetic\\
\midrule
5 & $12$--$13$ & 50 & $(10,10)$ & $(3,3,3,4)$ & 128.82 & 132.17\\
7 & $15$--$19$ & 70 & $(15,10)$ & $(4,4,6,6)$ & 171.69 & 175.04\\
9 & $22$--$23$ & 90 & $(19,14)$ & $(5,6,6,6)$ & 214.12 & 217.48\\
\bottomrule
\end{tabular}
\end{table}

Even the arithmetic model's column counts imply enormous storage.  An
information-theoretically packed vector of $M_s(\mathbf d)$ elements of
$\F_{19^4}$ requires
$M_s(\mathbf d)\log_2(19^4)$ bits.  Thus one vector---not a matrix, rank
profile, or solver state---already has the lower bounds in
\Cref{tab:vector-storage}; one sequential pass must move at least the same
amount of data.  These figures do not rule out streaming or structured
algorithms, but they prevent interpreting the arithmetic schedule as a
wall-clock or complete-circuit estimate.

\begin{table}[tbp]
\centering
\caption{Information-theoretic storage and one-pass traffic lower bounds for
one extension-field vector at the selected special-XL column count.  Decimal
GB/PB/EB units are used.}
\label{tab:vector-storage}
\small
\begin{tabular}{@{}c r r@{}}
\toprule
Level & $M_s(\mathbf d)$ & one packed vector / one full pass\\
\midrule
I   & $23{,}417{,}049{,}656$          & $49.74$ GB\\
III & $44{,}237{,}557{,}981{,}725{,}696$ & $93.96$ PB\\
V   & $236{,}094{,}291{,}515{,}544{,}000{,}000$ & $501.46$ EB\\
\bottomrule
\end{tabular}
\end{table}

The exhaustive search certificate for the optima in \Cref{tab:xl-optima} is
given in \Cref{app:cost-certificates}.  Under the $2^{32}$-gate envelope per retained target from
\Cref{sec:target-sampling}, generating the approximately $19^\tau$ structured
target trials costs less than $2^{74.48}$ gates for $\tau=10$ and
$2^{91.48}$ gates for $\tau=14$.  These additive terms do not change the
displayed hundredths.

Under the same per-trial Macaulay model, the AXN schedules remain below the
nominal references for rank defects $\Delta\le2,7,12$ at Levels I, III, and
V.  These thresholds control root-containing target--slice retries only; they
neither price a complete solver circuit nor certify unique recovery.  The
latter also requires $\bar\eta<1$ and the modeled one-dimensional Macaulay
kernel.

\subsection{The \texorpdfstring{$\ell=2$}{l=2} attacks}
\label{sec:l2-costs}

All three headline rows use the fixed-skew lift of
\Cref{prop:fixed-skew}.  The public preflight gives quotient rank $K$, affine
rank $M-K$, and the three residual systems listed in
\Cref{tab:l2-costs}.  CFXL specializes variables before XL-style linear
algebra, and Hashimoto refines this optimization for underdetermined
MQ~\cite{Beullens2024SNOVA,Hashimoto2023}.  In Hashimoto's objective, $a$ is
the square-subsystem size and $k$ is a second specialization parameter.  The
complete objective, boundary conventions, and exact minimizers are recorded
in \Cref{app:generic-mq-certificate}.

Every recovered fixed-skew root passes the block-symmetry rejection test.  For
a global polar-rank lower bound $r_*$, target retries add at most
\[
 \log_2\!\left(1+(1-19^{-K})19^{K-r_*}\right)
\]
bits to the expected-cost exponent; $r_*\ge2K$ removes this retry by making
every target solvable.  The first three rows of \Cref{tab:l2-costs} give the
resulting selected per-solve schedules and the corresponding global-rank
thresholds.

\begin{table}[tbp]
\centering
\caption{$\ell=2$ systems and global polar-rank thresholds.  Model is the
pinned published per-solve operation model; AXN uses
$g_1^{\mathrm{AXN}}=441$.  For fixed-skew, ``For gap'' is the smallest global
rank for which the root-availability retry bound keeps expected AXN arithmetic
below the nominal NIST reference.  For that construction,
``all-target'' makes every target solvable and every recovered root passes the
block-symmetry rejection test.  For zero-offset it proves only that an accepted root
exists; the row is not selected as a headline attack.}
\label{tab:l2-costs}
\scriptsize
\setlength{\tabcolsep}{3.25pt}
\begin{tabular}{@{}l l r r r r r@{}}
\toprule
Parameters & mode & eqs./vars. & model & AXN & \shortstack{for\\gap} & \shortstack{all-target\\rank}\\
\midrule
$(48,16,2,2)$ & fixed-skew & $48/112$ & 129.44 & 133.75 & 46  & 96\\
$(72,24,2,2)$ & fixed-skew & $72/168$ & 183.44 & 187.74 & 68  & 144\\
$(96,32,2,2)$ & fixed-skew & $96/224$ & 237.84 & 242.15 & 89  & 192\\
$(96,32,2,2)$ & zero-offset per solve & $96/256$ & 235.19 & 239.50 & -- & 240\\
\bottomrule
\end{tabular}
\end{table}

The zero-offset row is a per-solve sensitivity only.  Its dash records that no
polar-rank threshold controls which root the priced solver returns, so rank
alone does not imply an expected accepted-root cost.  Its all-target threshold
includes block acceptance through \eqref{eq:l2-rejection}; for fixed-skew the
condition is simply $r_*\ge2K$.  The complete zero-offset retry, target,
and solver-output analysis is in \Cref{app:zero-offset}.

\section{Implementation and Evidence}\label{sec:evidence}

We report an implementation of the public key generator, verifier-side feature
maps, symmetry quotient, affine reduction, structured lift, rank tests, and
cost searches over $\F_{19}$.  Machine-readable certificates and supporting
artifacts exist for these computations, but they are not bundled with this
manuscript.  The computations are therefore labeled
\emph{author-reported evidence}, not evidence independently reproduced from an
attached archive; \Cref{app:artifact} states
the release contract needed to change that status.  The reported implementation does \emph{not} include a production-size
end-to-end forgery, a production-size Macaulay solve, or a wall-clock attack
timing.  Its full-size experiments concern public preprocessing and rank; its
solver experiments use smaller tractable profiles.  The checks therefore have
different logical force.  The table
below is the intended reading guide for every experimental statement in this
section.

\begin{table}[H]
\centering
\caption{Evidence hierarchy.  ``Exact'' here means exact arithmetic for the
reported concrete object, not independent reproduction or an exhaustive
theorem about all production directions.}
\label{tab:evidence-hierarchy}
\small
\renewcommand{\arraystretch}{1.12}
\begin{tabularx}{0.96\textwidth}{@{}>{\raggedright\arraybackslash\bfseries}p{0.24\textwidth}>{\raggedright\arraybackslash}p{0.30\textwidth}>{\raggedright\arraybackslash}X@{}}
\toprule
Check & What is computed & What it does---and does not---establish\\
\midrule
Public preflight & Exact ranks and rejection conditions for the chosen key,
relation, offsets, and slice & The exact reduction is valid for that concrete
attack instance.\\
Rank campaign & Exact ranks for many selected nonzero output or direction
vectors & Finite evidence about the unenumerated global minimum rank; not a
certificate for all directions.\\
Macaulay experiment & Exact kernels at smaller tractable dimensions & Evidence
for the production solver model; not a production timing or rank proof.\\
Circuit validation & Exhaustive checking and comparison with a separately
implemented evaluator for the reported field primitive & Exact cost and
functionality of that primitive, but not of a
complete solver.\\
\bottomrule
\end{tabularx}
\end{table}

Projective directions were defined in \Cref{sec:rank-primer}.  Sampling
many distinct directions can expose low-rank structure, but only exhaustive
enumeration or a proof can certify the global minimum required by the
headline all-target-root conditions.

\subsection{Public preflight and verifier regression}

For a structured $\ell=4$ instance, the preflight checks
\[
 \rank Q_R=K,\qquad \rank C_{R,\mathbf v}=M-K,\qquad
 \rank\mathcal O_{R,\mathbf v}=m_1\ell^2,
\]
together with $\mathcal F_{R,\mathbf v}H_{R,\mathbf v}=0$, the predicted $A$-kernel
dimension, $A$-invariance, the restricted quadratic rank, and rejection
compatibility.  In the reported campaign, all 121 regenerated full-size
instances pass,
including the official Level-I known-answer key.  The first deterministic
offset candidate passes for 114 instances; the second passes for the remaining
seven.

For fixed-skew, the preflight checks the relation and rejection identity and
the quotient and affine ranks, including the ranks after adjoining the
relevant public constant column.  All 60 full-size keys tested in the reported
campaign pass.  The
reported regressions also reproduce the official
Level-I public key byte for byte,
verifying its known-answer signature against the official hash target, and
checking the fixed output-order permutation described in
\Cref{sec:affine-columns}.  Eighteen additional regressions cover all nine
Version~2.3 shapes.  A separately implemented Gauss--Jordan routine agrees with
the primary rank routine on 2,036 reported matrices, including 36
production-sized matrices.

\subsection{Rank evidence}

\paragraph{Structured slices.}
Every selected slice is tested on all coefficient-basis directions (one
nonzero output coefficient at a time) and 1,024 deterministic dense directions
(many nonzero coefficients).  Selected slices also receive complete normalized
weight-two tests (two nonzero coefficients, modulo overall scaling) or
additional dense tests.  All 102,796 tested
cross-polar maps have maximal rank $h=4s$; the per-shape counts are in
\Cref{tab:cross-polar}.  The dual campaign evaluates 21,992 directional maps
$T_y$ on the same 13 slices, and every one is surjective of rank $K$.
All 124,788 coefficient and direction vectors used in these two campaigns are
projectively distinct within their respective slice records.

\begin{table}[tbp]
\centering
\caption{Author-reported exact cross-polar ranks for the tested structured
$\ell=4$ slice--direction pairs.}
\label{tab:cross-polar}
\small
\setlength{\tabcolsep}{5.2pt}
\begin{tabular}{@{}l r r r r@{}}
\toprule
Parameters & slices & $h=4s$ & rank tests & sampled minimum\\
\midrule
$(28,5,4,4)$  & 3 & 40 & 26,296 & 40\\
$(28,4,4,5)$  & 2 & 40 & 24,198 & 40\\
$(40,7,4,4)$  & 2 & 60 & 45,658 & 60\\
$(38,5,4,5)$  & 2 & 60 & 2,188  & 60\\
$(50,9,4,4)$  & 2 & 76 & 2,228  & 76\\
$(52,6,4,6)$  & 2 & 76 & 2,228  & 76\\
\midrule
Total          & 13 & -- & 102,796 & maximal\\
\bottomrule
\end{tabular}
\end{table}

\paragraph{Fixed-skew systems.}
A broad campaign tests every coefficient-basis direction together with
combinations having two, three, or many nonzero output coefficients.  A second campaign targets
the 21 exceptional projective directions that make the following block-local
polar operator have rank two rather than four.  For an output direction $b$
and a signature block $c$, let
\[
 \Lambda_{b,c}:\Mat_2(\F_{19})\longrightarrow\Mat_2(\F_{19})^*,
 \qquad
 \Lambda_{b,c}(X)(Y)=b\cdot\beta_c(X,Y),
\]
where $\beta_c$ is the restriction of the residual polar map to variables in
block $c$.  The campaign tests all 381 local
projective directions in one block, every exceptional direction in every
remaining block, and deterministic two-block exceptional combinations.
Together the campaigns perform 3,888 exact rank evaluations; their sampled
minima and the rigorous thresholds used by the cost analysis appear in
\Cref{tab:fixed-skew-evidence}.

\begin{table}[tbp]
\centering
\caption{Author-reported full-size polar-rank evidence for the fixed-skew
$\ell=2$ systems.
``For nominal gap'' is the minimum global polar rank that preserves a positive
AXN nominal arithmetic gap in \Cref{tab:l2-costs}; ``all-target root'' guarantees root
existence for every target, not success of a particular solver run.}
\label{tab:fixed-skew-evidence}
\small
\setlength{\tabcolsep}{4.5pt}
\begin{tabular}{@{}l r r r r r@{}}
\toprule
Parameters & preflights & rank tests & sampled min. & \shortstack{for nominal\\gap} & all-target root\\
\midrule
$(48,16,2,2)$ & 20 & 1,104 & 110 & 46 & 96\\
$(72,24,2,2)$ & 20 & 1,296 & 167 & 68 & 144\\
$(96,32,2,2)$ & 20 & 1,488 & 223 & 89 & 192\\
\midrule
Total & 60 & 3,888 & -- & -- & --\\
\bottomrule
\end{tabular}
\end{table}

The separate zero-offset rank campaign and its solver-output sensitivity are
reported with that unselected construction in \Cref{app:zero-offset}.

\subsection{Solver and probability checks}

The reported special-XL search examines all 42,356 admissible multidegrees at modeled
work no greater than the incumbent and recovers the three optima in
\Cref{tab:xl-optima}.  At tractable sizes, the reported campaign contains 128 four-block
Macaulay matrices: 64 random instances have full column rank, while 64 planted
instances, each generated with a known root, have a one-dimensional right
kernel containing the planted evaluation vector.  Because the columns include
the constant and degree-one monomials, that one-dimensional kernel also rules
out a second base-field root in each tested instance.  The campaign sampled
planted instances rather than first conditioning on $Z=1$, so it is supporting
evidence, not a certified lower bound for $p_{\rm ker}^{(1)}$.  A 48-matrix two-block control
contains one planted case with a two-dimensional right kernel, so the
corresponding $\ell=2$ special-XL estimates are not used.  These are
small-profile generic-system experiments, not production SNOVA timings.

Exhaustive checks over $\F_3$ and $\F_5$ verify the exact first and factorial
second moments, the compatible-collision formula, the rank-spectrum upper bound,
and the singleton and planted-distribution inequalities on 110 fully enumerated
quadratic maps.  Seventeen maps have nonmaximal cross-polar rank; in twelve,
compatibility makes $\eta_0$ strictly smaller than $\bar\eta$.  Finally, on
six actual cost-driving $\ell=2$ cores, the homogeneous quadratic outputs are
linearly independent and the common polar radical---the set of directions
annihilated by every polar form---is zero.  These are genericity diagnostics,
not solving-degree measurements.

\subsection{Scope of the estimates}

The evidence hierarchy in \Cref{tab:evidence-hierarchy} separates exact public
preflights from finite rank campaigns and smaller solver experiments.  The
remaining production-scale inputs are the global polar spectra and solver
behavior; no production-size end-to-end solve is claimed.

\paragraph{What would turn the estimates into end-to-end certificates.}
The rank assumptions would require either an algebraic proof of the global
minimum over every nonzero output direction or a complete enumeration enabled
by additional structure.  The solver assumptions would require
production-size Macaulay or Gr\"obner runs, or a validated implementation
whose observed ranks and costs replace the semi-regular models.  A complete
gate comparison would additionally need to price control flow, pivoting,
indexing, memory, and data movement.

\section{Countermeasures}\label{sec:countermeasures}

The design lesson is that security must be based on distinct quadratic
polynomials, not on the number of ordered labels used to name them.  A security
analysis for any symmetric redesign should compute the actual quotient rank
$\rho_{\rm quad}$ and use the unconditional cap
\begin{equation}\label{eq:countermeasure-K}
 \rho_{\rm quad}\le K_d
 =\min\left\{M,m_1\binom{d+1}{2}\right\},
 \qquad d=\dim_{\F_q}\F_q[S].
\end{equation}
Equality with $K_d$ is a public rank condition, not a consequence of symmetry
alone.  Outer randomization does not restore the alternating directions removed
before the outer map.

Candidate design responses include increasing parameters after accounting for
\eqref{eq:countermeasure-K},
removing public-matrix symmetry, or redesigning the left and right power actions so
that exchanging their labels changes the quadratic polynomial.  None is a
drop-in security proof: each changes other algebraic structure and requires
fresh cryptanalysis, performance evaluation, and parameter selection.  The
block-symmetry rejection test is not sufficient: the fixed-skew $\ell=2$
lift makes every block nonsymmetric by construction, and the structured
$\ell=4$ preflight can certify the same property block by block.  Returning
unchanged to the submitted $q=16$ design would instead restore the known
characteristic-two attack surface.

A revised security analysis should publish the quotient rank, affine
target-preselection rank, stable-kernel rank when applicable, the polar or cross-polar condition
used for target availability, and solver logs for the actual cost-driving
system.

\section{Conclusion}\label{sec:conclusion}

The central vulnerability is an exact duplicate-feature identity.  \SNOVA{}
labels a quadratic power-pair feature by an ordered pair $(\xi,\eta)$, but
symmetry of $S$ and of the public base matrices makes that polynomial identical to the
one labeled $(\eta,\xi)$.  The public $q=19$ verifier therefore has at most
$m_1\binom{\ell+1}{2}$ distinct power-pair features where the ordered
representation has $m_1\ell^2$.  Because the identification occurs before the
public outer map, no later randomization can restore the missing directions.

The rest of the paper turns that identity into an exact reduction of all $M$
equations after imposing the attacker-chosen column relation.  A stable
$A$-linear slice preserves
the four-block structure needed for the six $\ell=4$ special-XL systems.
Fixed-skew relations give the selected reductions for all three $\ell=2$
systems while satisfying the block-symmetry rejection test; a zero-offset construction is
retained only as a solver-output-sensitive alternative.  Exact moment formulas distinguish root
existence, unique-root retries, and the separate question of whether the
chosen Macaulay matrix exposes the root.  Every candidate is checked by the
unmodified verifier.

Under the stated global-rank and production-solver assumptions, the modeled
expected operation counts are \LoneModelRange, \LthreeModelRange, and
\LfiveModelRange.  Replacing the field primitive charged by those models with
author-reported signed AXN circuits, as described in
\Cref{app:field-circuits}, gives
\LoneAXNRange, \LthreeAXNRange, and
\LfiveAXNRange.  These arithmetic-only figures leave
\AXNNominalGapRange{} bits below NIST's nominal gate references and
\AXNExpectedAESGapRange{} bits below expected AES key search, or
\AXNFullAESGapRange{} bits below full-key-space AES enumeration, using the
reported calibration circuits in the same basis.  The comparison remains
resource-asymmetric because the MQ schedules omit substantial implementation
costs.

The exact and conditional claims remain distinct.  The quotient and reduction
theorems apply to every key satisfying their public preflights, whereas uniform
retry sampling additionally uses the stated random-oracle experiment and the
headline exponents additionally use uncertified global-rank and production
solver models.  The circuit conversion prices dominant arithmetic rather than
a complete solver, so an end-to-end security classification still requires
global rank certificates and production-scale measurements.

For a redesign, the immediate lesson is simple: count distinct quadratic
polynomials after every public symmetry, not the labels used to name them.

\appendix
\setcounter{section}{0}
\renewcommand{\thesection}{\Alph{section}}
\renewcommand{\theHsection}{app.\Alph{section}}


\section{Deferred Structural Proofs}\label{app:deferred-proofs}

This appendix supplies the longer proofs for the stable-kernel and
stable-slice results used in the main construction.

\subsection{Stable-lift kernel}\label{app:proof-stable-kernel}

\begin{proof}[Proof of \Cref{thm:stable-kernel}]
Choose an $\F_q$-basis $\{\xi_1,\ldots,\xi_d\}$ of $A$.  By bilinearity, the
$m_1d^2$ basis-indexed rows
$w^\transpose\Sigma_{\xi_a}P_i\Sigma_{\xi_b}$ span the family in
\eqref{eq:Uw}; some of these rows may coincide or vanish.  This proves the
dimension bound without asserting that the all-$\xi,\eta$ family has that
literal cardinality.  For
$\zeta\in A$,
\[
 \bigl(w^\transpose\Sigma_\xi P_i\Sigma_\eta\bigr)\Sigma_\zeta
 =w^\transpose\Sigma_\xi P_i\Sigma_{\eta\zeta},
\]
so $\mathcal U_w$ is closed under right multiplication by $A$.

Consider a cross term produced by substituting
$u_j=\Sigma_{R_j}u+\Sigma_{\Gamma_j}w$ into a feature.  A term with $w$ on the left
already has row coefficient
$w^\transpose\Sigma_\xi P_i\Sigma_\eta$.  A term with $w$ on the right can be
transposed; using $S^\transpose=S$, $P_i^\transpose=P_i$, and commutativity of
$A$, its row coefficient has the same form with the two algebra labels
reversed.  Thus every row of $\mathcal F_{R,\mathbf v}$ lies in $\mathcal U_w$.
Consequently every row of
$C_{R,\mathbf v}=W_R\mathcal E\mathcal F_{R,\mathbf v}$ lies there as well.  Right-$A$
stability then puts every row of the orbit matrix $\mathcal O_{R,\mathbf v}$ in
$\mathcal U_w$.

If $\rank\mathcal O_{R,\mathbf v}=m_1d^2$, then
\[
 m_1d^2=\dim\operatorname{row}(\mathcal O_{R,\mathbf v})
 \le\dim\mathcal U_w\le m_1d^2,
\]
so $\operatorname{row}(\mathcal O_{R,\mathbf v})=\mathcal U_w$.  Since the rows of
$\mathcal F_{R,\mathbf v}$ belong to this row space,
$\mathcal F_{R,\mathbf v}\ker\mathcal O_{R,\mathbf v}=0$.  To see $A$-stability of the
kernel, let $x\in\ker\mathcal O_{R,\mathbf v}$ and $\zeta\in A$.  Every product
$S^a\zeta$ is an $\F_q$-linear combination of
$1,S,\ldots,S^{d-1}$, so
$C_{R,\mathbf v}\Sigma_{S^a}\Sigma_\zeta x=0$ for all $a$; hence
$\Sigma_\zeta x\in\ker\mathcal O_{R,\mathbf v}$.  Rank--nullity gives
$\dim_{\F_q}H_{R,\mathbf v}=N_0-m_1d^2$.  When $A\cong\F_{q^\ell}$, $H_{R,\mathbf v}$ is an
$A$-vector space, and division by $\ell$ gives
$\dim_A H_{R,\mathbf v}=n-m_1\ell$.
\end{proof}

\subsection{Stable affine slices}\label{app:proof-stable-slices}

\begin{proof}[Proof of \Cref{cor:stable-slices}]
The first two ranks in \eqref{eq:stable-preflight} give the full-rank case
$\rho=K$, $\lambda=M-K$ of \Cref{prop:affine-reduction}; in particular, every target determines a
nonempty affine space $a(t)+\ker C_{R,\mathbf v}$.  By its definition in
\eqref{eq:orbit-matrix}, the orbit matrix contains $C_{R,\mathbf v}$ as its
first block, so
$H_{R,\mathbf v}\subseteq\ker C_{R,\mathbf v}$.  Hence
$a(t)+y+L\subseteq a(t)+\ker C_{R,\mathbf v}$ whenever
$y\in\ker C_{R,\mathbf v}$.  Because
$L\subseteq H_{R,\mathbf v}$, \Cref{thm:stable-kernel} gives
$\mathcal F_{R,\mathbf v}L=0$, and therefore $\mathcal F_{R,\mathbf v}u$ is constant on that
slice.

By \Cref{lem:spectral-quotient}, the unordered power-pair features can be
changed invertibly to eigenblock coordinates.  Because $Q_R$ has full column
rank, row reduction of the complete output system isolates $K$ rows whose
homogeneous quadratic parts are an invertible linear image of
$z_{\rm sym}$.  Applying the inverse linear change of feature coordinates to
those rows expresses the
quadratic parts in eigenblock coordinates.  Full rank of the quadratic feature
map restricted to $L$ ensures that this restriction loses no eigenblock
component.  Translating a diagonal feature introduces only
quadratic, linear, and constant terms from one block; translating an
off-diagonal feature introduces terms from its two blocks.  If $h_a$ is a
homogenizer for block $a$, multiplication by the unique powers of the $h_a$
turns these affine polynomials into exact multidegrees $2\mathbf e_a$ and
$\mathbf e_a+\mathbf e_b$.

Distinct unordered multidegrees have disjoint monomial supports.  Each
component is generated by the $m_1$ base forms and therefore has rank at most
$m_1$.  The total restricted rank is the sum of these component ranks.  Since
there are $\binom{\ell+1}{2}$ components and the total is
$K=m_1\binom{\ell+1}{2}$, every component must have rank exactly $m_1$.
\end{proof}

\section{The Zero-Offset \texorpdfstring{$\ell=2$}{l=2} Alternative}
\label{app:zero-offset}

Set $R_0=I$, $R_1=0$, and use zero offsets.  Each signature block then has
the form
\[
 \begin{pmatrix}x&0\\ y&0\end{pmatrix},
\]
so it is symmetric exactly when $y=0$.  The $n$ symmetry conditions are
nonzero hyperplanes on disjoint block coordinates.  A relation that instead
forced the two off-diagonal entries to agree would make every block symmetric
and would be unusable.

Let $N=N_0$ be the residual-domain dimension and let
$g:\F_q^N\to\F_q^K$ be the residual map of degree at most two.  For each
nonzero $b\in\F_q^K$, let $r(b)$ be the polar rank of $b\cdot g$ and put
$r_*=\min_{b\ne0}r(b)$.  Define
\[
 t_{\rm rej}=\lfloor n/4\rfloor+1,
 \qquad B=\binom n{t_{\rm rej}},
 \qquad \alpha=1-q^{-K},
\]
and let $\mathcal X_{\rm acc}$ be the assignments accepted by the
block-symmetry rejection test.  Since at least $t_{\rm rej}$ symmetric blocks
are needed for rejection, the exact density of accepted assignments is
\begin{equation}\label{eq:l2-accepted-density}
 a_n\coloneqq\frac{|\mathcal X_{\rm acc}|}{q^N}
 =q^{-n}\sum_{k=0}^{\lfloor n/4\rfloor}\binom nk(q-1)^{n-k}.
\end{equation}

\begin{proposition}[Accepted roots for the zero-offset $\ell=2$ lift]\label{prop:l2-rejection}
Every target has an accepted root if
\begin{equation}\label{eq:l2-rejection}
 1-Bq^{-t_{\rm rej}}>\alpha(1+B)q^{K-r_*/2}.
\end{equation}
For a uniform target $z\in\F_q^K$,
\begin{equation}\label{eq:l2-accepted-target}
 \Pr\bigl[g^{-1}(z)\cap\mathcal X_{\rm acc}\ne\varnothing\bigr]
 \ge \frac{a_n}{1+\alpha q^{K-r_*/2}}.
\end{equation}
Moreover, if $1-\alpha q^{K-r_*/2}>0$, every target $z$ satisfies
\begin{equation}\label{eq:l2-accepted-fraction}
 \frac{|g^{-1}(z)\cap\mathcal X_{\rm acc}|}{|g^{-1}(z)|}
 \ge
 1-\frac{Bq^{-t_{\rm rej}}+\alpha Bq^{K-r_*/2}}
          {1-\alpha q^{K-r_*/2}},
\end{equation}
whenever the right-hand side is positive.
\end{proposition}

We first record the character-sum estimate used in the proof.

\begin{lemma}[Quadratic fiber bound]\label{lem:quadratic-fiber-bound}
Let $q$ be odd and let $g:\F_q^N\to\F_q^K$ be quadratic.  For each
$b\in\F_q^K\setminus\{0\}$, let $r(b)$ be the polar rank of the scalar
quadratic $b\cdot g$, put $r_*=\min_{b\ne0}r(b)$, and set
$\alpha=1-q^{-K}$.  Then every target $z\in\F_q^K$ satisfies
\begin{equation}\label{eq:quadratic-fiber-bound}
 \left|\,|g^{-1}(z)|-q^{N-K}\right|
 \le \alpha q^{N-r_*/2}.
\end{equation}
More generally, if $W\subseteq\F_q^N$ is an affine subspace of codimension
$c$, then
\begin{equation}\label{eq:quadratic-affine-upper}
 |\{x\in W:g(x)=z\}|
 \le q^{N-c-K}+\alpha q^{N-r_*/2}.
\end{equation}
\end{lemma}

\begin{proof}
Fix a nontrivial additive character $\psi$ of $\F_q$.  Fourier inversion gives
\[
 |g^{-1}(z)|
 =q^{-K}\sum_{b\in\F_q^K}\psi(-b\cdot z)
   \sum_{x\in\F_q^N}\psi(b\cdot g(x)).
\]
The $b=0$ term is $q^{N-K}$.  For $b\ne0$, let $r=r(b)$.  An invertible
linear change of variables splits the polar form into a nondegenerate
$r$-dimensional part and a radical of dimension $N-r$.  If the linear part of
$b\cdot g$ is nonzero on the radical, summing over the radical gives zero.
Otherwise the radical contributes $q^{N-r}$, while diagonalizing the
nondegenerate quadratic part reduces its sum to a product of one-dimensional
Gauss sums of absolute value $q^{1/2}$.  In either case the inner sum has
absolute value at most $q^{N-r/2}$~\cite{Lam2005}.  Summing the $q^K-1$
nonzero Fourier coefficients proves \eqref{eq:quadratic-fiber-bound}.

Write $W=a+U$, where $U$ has codimension $c$.  Translation does not change a
quadratic polynomial's polar form.  Restricting a bilinear form to
$U\times U$ lowers its rank by at most $2c$: after extending a basis of $U$ to
one of $\F_q^N$, deleting the last $c$ rows and columns changes matrix rank by
at most $2c$.  Thus every nonzero scalar combination of $g|_W$ has polar rank
at least $r_*-2c$.  If $r_*\ge2c$, the same Fourier calculation on the
$(N-c)$-dimensional affine space $W$ bounds each nonzero character sum by
\[
 q^{(N-c)-(r_*-2c)/2}=q^{N-r_*/2}.
\]
If $r_*<2c$, the trivial estimate gives
$|\{x\in W:g(x)=z\}|\le q^{N-c}$.  In this branch,
\[
 \frac{q^{N-c-K}+\alpha q^{N-r_*/2}}{q^{N-c}}
 =q^{-K}+(1-q^{-K})q^{c-r_*/2}\ge1,
\]
because $c-r_*/2>0$.  Hence the right-hand side of
\eqref{eq:quadratic-affine-upper} still dominates the trivial estimate.  This
proves \eqref{eq:quadratic-affine-upper} in both cases.
\end{proof}

\begin{proof}[Proof of \Cref{prop:l2-rejection}]
By \Cref{lem:quadratic-fiber-bound}, every target fiber has between
$q^{N-K}-\alpha q^{N-r_*/2}$ and
$q^{N-K}+\alpha q^{N-r_*/2}$ roots.  A rejected root lies on
the intersection of some $t_{\rm rej}$ block-symmetry hyperplanes.  Each such
intersection is an affine subspace of codimension $t_{\rm rej}$, so
\eqref{eq:quadratic-affine-upper} bounds the number of roots it contains by
$q^{N-t_{\rm rej}-K}+\alpha q^{N-r_*/2}$.  Subtracting the union bound
over the $B$ intersections from the full-fiber lower bound gives
\[
 q^{N-K}(1-Bq^{-t_{\rm rej}})
 -\alpha q^{N-r_*/2}(1+B).
\]
This is positive exactly under the sufficient inequality
\eqref{eq:l2-rejection}.

Every fiber has size at most
$q^{N-K}(1+\alpha q^{K-r_*/2})$.  Since
$|\mathcal X_{\rm acc}|=a_n q^N$, double-counting accepted assignments by their images
proves \eqref{eq:l2-accepted-target}.  Finally, let
$T_z=|g^{-1}(z)|$ and let $R_z$ be the number of rejected roots.  The accepted
fraction is $1-R_z/T_z$.  Applying the union-bound upper estimate to $R_z$ and
the character-sum lower estimate to $T_z$ gives
\eqref{eq:l2-accepted-fraction}.
\end{proof}

The three conclusions answer different questions.  For the Level-V system,
the global condition $r_*\ge240$ implies \eqref{eq:l2-rejection} and hence an
accepted root for every target.  Inequality \eqref{eq:l2-accepted-target}
controls the fraction of uniform targets having an accepted root, whereas
\eqref{eq:l2-accepted-fraction} controls the accepted population fraction
inside each target fiber.  None controls which root a deterministic or
nonuniform solver returns.

\paragraph{Target-generation budget.}
The Level-V alternative needs on average $19^{32}$ retained unbiased targets
before its $M-K=32$ affine consistency coordinates take their prescribed
values.  The worst-case chunk-acceptance bound in
\Cref{sec:target-sampling} therefore gives at most
$6.6\cdot19^{32}$ raw message--salt candidates on average.  If
$p_{\rm chunk}$ is the actual chunk-acceptance probability, exhausting one
128-bit salt space yields
\[
 p_{\rm chunk}\frac{2^{128}}{19^{32}}
 \ge 0.00062357\ldots
\]
expected candidates passing both rejection conversion and affine
preselection.  Under independent random-oracle samples, the probability of at
least one such candidate in a complete salt space exceeds $0.0623\%$, so the
expected number of chosen-message prefixes is below $1{,}605$.  The prefix
digests are checked for uniqueness as described in
\Cref{sec:target-sampling}, and the forged message is kept distinct from every
signing-query message.  This is an existential chosen-message forgery, not a
signature on an externally fixed message.

Under the deliberately loose envelope of fewer than $2^{32}$ gates per
retained target, including the unbiased chunk sampler, target preselection
costs less than
\[
 2^{32}19^{32}=2^{167.93368}.
\]
Even charging the same envelope to every raw candidate gives at most
$2^{170.64716}$.  Both are additive costs more than 68 exponent bits below
the zero-offset AXN generic-MQ schedule; neither prices the separate task of
obtaining an accepted root from that solver.

\paragraph{Solver-output sensitivity.}
For Level V, target preselection leaves 96 equations in all 256 variables.
The pinned generic-MQ objective gives per-solve Model and AXN exponents
$235.19474$ and $239.50086$.  Target generation occurs before a solve and is
therefore additive, not multiplicative.  If the global polar-rank minimum
satisfies $r_*=240$, \eqref{eq:l2-accepted-fraction} proves only
\[
 p_{\rm acc}\ge0.1224162627718736
\]
for the accepted population fraction in each target fiber.  If one
\emph{additionally} assumes that each solve returns an independent uniform
root of that fiber, accepted-root recovery takes an expected
$1/p_{\rm acc}=8.16884928=2^{3.03013286}$ solves.  Charging that factor gives
Model and AXN exponents $238.22487$ and $242.53100$, both slightly worse than
fixed-skew.  The pinned solver is not proved to sample roots uniformly, so
these are sensitivity figures rather than an attack-cost theorem.

The reported zero-offset rank campaign performs 29,424 exact evaluations on
two keys per parameter shape and finds sampled minima 127, 190, and 254 in
dimensions 128, 192, and 256.  If the Level-V global minimum were at least
254, \eqref{eq:l2-accepted-fraction} would give an accepted fraction exceeding
$0.999999999$ in every target fiber.  The finite campaign supports but does
not certify that hypothesis, and the solver-output distinction remains.

\section{Deferred Probability Proofs}\label{app:probability-proofs}

The next three proofs establish the root-collision, planted-distribution, and
full-space statements used in \Cref{sec:slices}.

\subsection{Root collisions and the rank-spectrum bound}\label{app:proof-coset-success}

\begin{proof}[Proof of \Cref{thm:coset-success}]
Average over the $q^{\dim V-h}$ cosets and $q^K$ targets.  Every point of $V$
contributes once to the first moment, which gives \eqref{eq:EZ}.  For an
ordered pair of distinct points in one coset, write the second point as $x+y$
with $y\in L\setminus\{0\}$.  The collision equation is
\[
 g(x+y)=g(x)
 \quad\Longleftrightarrow\quad
 T_y(x)=-(g(y)-g(0)).
\]
It has $\kappa(y)q^{\dim V-r(y)}$ solutions.  Dividing the total number of
ordered collisions by the number of target--coset pairs proves
\eqref{eq:factorial-moment}.

For the rank-spectrum identity, note that
$q^{-r(y)}=q^{-K}|\ker T_y^*|$ and double-count pairs $(b,y)$ satisfying
$b\cdot T_y=0$.  The term $b=0$ contributes $q^h-1$.  For $b\ne0$, the number
of nonzero $y\in L$ annihilated by $b\cdot T_y$ is
$q^{h-e_L(b)}-1$.  Therefore
\[
 \bar\eta
 =q^{-K}\left((q^h-1)+\sum_{b\ne0}(q^{h-e_L(b)}-1)\right)
 =\mu(1+\Theta_L(g))-1.
\]
The inequality $\eta_0\le\bar\eta$ is immediate.

Using $\mathbb E[Z^2]=\mu(1+\eta_0)$ gives
\[
 \operatorname{Var}(Z)
 =\mu(1+\eta_0-\mu)
 \le\mu^2\Theta_L(g),
\]
and the second-moment inequality proves \eqref{eq:coset-success}.  Finally,
$e_*\ge h-\Delta$ implies
$\Theta_L(g)\le(q^K-1)q^{-h+\Delta}<q^{\tau+\Delta}$.
\end{proof}

\subsection{Unique-root trials and the planted distribution}\label{app:proof-singleton-coset}

\begin{proof}[Proof of \Cref{cor:singleton-coset}]
For every nonnegative integer $z$,
\[
 \mathbf 1_{z=1}\ge z-z(z-1),\qquad
 \mathbf 1_{z>0}\ge z-\binom z2,\qquad
 \mathbf 1_{z\ge2}\le\binom z2.
\]
Apply these inequalities to $Z$ and use
\eqref{eq:factorial-moment} together with $\eta_0\le\bar\eta$.
Dividing the upper bound on $\Pr[Z\ge2]$ by
\eqref{eq:bonferroni-success} proves \eqref{eq:conditional-multiple}.

The planted law size-biases a target--coset pair by $Z$.  Hence
\[
 \Pr_{\mathsf P}[Z\ge2]
 =\frac{\mathbb E[Z\mathbf 1_{Z\ge2}]}{\mu}
 \le\frac{\mathbb E[Z(Z-1)]}{\mu}=\eta_0.
\]
Conditioned on $Z=1$, both $\mathsf U$ and $\mathsf P$ are uniform on the same
set of target--coset pairs.  If two distributions have the same conditional
law on an event, their total variation distance is at most the larger mass they
assign to its complement.  Here the complement is $\{Z\ge2\}$.  The same indicator inequality gives the sharper exact-parameter lower bound
$\Pr[Z>0]\ge\mu(1-\eta_0/2)$.  Combining it with
$\Pr[Z\ge2]\le\mathbb E[Z(Z-1)]/2=\mu\eta_0/2$ gives
\[
 \Pr_{\mathsf U}[Z\ge2]
 \le \frac{\eta_0}{2-\eta_0}\le\eta_0
 \qquad(\eta_0<1).
\]
The displayed size-bias bound gives
$\Pr_{\mathsf P}[Z\ge2]\le\eta_0$.  Both are at most $\bar\eta$, proving the
total-variation claim.
\end{proof}

\subsection{Full-space target availability}\label{app:proof-full-space}

\begin{proof}[Proof of \Cref{cor:full-space}]
Take $L=V$ in \Cref{thm:coset-success}.  Then $e_L(b)$ is the polar rank and
$\Theta_L(g)\le(q^K-1)q^{-r_*}=\alpha q^{K-r_*}$.  For the last statement, \eqref{eq:quadratic-fiber-bound} gives
\[
 |g^{-1}(z)|\ge q^{N-K}-\alpha q^{N-r_*/2}>0
\]
when $r_*\ge2K$.
\end{proof}

\section{Implementation Conventions and Circuit Calibration}
\label{app:implementation}

\subsection{Output-coordinate permutation}\label{app:output-permutation}

The reference public-output block $T_\mu$ has $r$ rows and $\ell$ columns.
For $0\le i<r$ and $0\le j<\ell$, the reference C verifier stores
\[
 (t_{\rm off})_{\mu r\ell+i\ell+j}=T_\mu[i,j].
\]
Our equations instead flatten the transposed $\ell\times r$ block row by row.
Thus the required permutation is
\[
  t_{\mu r\ell+jr+i}
  =(t_{\rm off})_{\mu r\ell+i\ell+j}.
\]
For a vector $y_{\rm off}$ in the official order, define
\[
 (\Pi_{\rm out}y_{\rm off})_{\mu r\ell+jr+i}
 =(y_{\rm off})_{\mu r\ell+i\ell+j}.
\]
Then $t=\Pi_{\rm out}t_{\rm off}$ and
$\mathcal E=\Pi_{\rm out}\mathcal E_{\rm off}$.  The indices above remain
valid when $r\ne\ell$; the transformation is only a coordinate permutation
and preserves ranks and solutions.

\subsection{Fixed-skew identity}\label{app:fixed-skew}

For the official $\ell=2$ matrix, direct arithmetic gives
\[
 9I+10S=\begin{pmatrix}0&1\\1&7\end{pmatrix},
 \qquad e_0^\transpose(9I+10S)=e_1^\transpose.
\]
Together with \eqref{eq:fixed-skew-offset}, this proves
\eqref{eq:fixed-skew-identity} without a rank or genericity assumption.  The
remaining public preflight checks $\rank Q_R=K$ and
$\rank C_{R,\mathbf v}=M-K$.

The quotient matrix depends only on the fixed public outer map and the
relation, not on a generated key.  Exhaustive
enumeration of all $19^2=361$ relations $R_1=aI+bS$ gives quotient rank $K$
for every relation on all three $\ell=2$ parameter shapes.  The resulting
residual dimensions, exact optimizer witnesses, converted costs, and
polar-rank thresholds are recorded in
\Cref{tab:l2-costs,tab:generic-mq-witnesses,tab:fixed-skew-evidence}.

\subsection{Field-arithmetic circuits}\label{app:field-circuits}

For $\F_{19}$, the reported circuits use canonical five-bit representatives
and exact Barrett reduction.  One signed multiply-accumulate, $c+ab$ or
$c-ab$, uses 402 or 441 AXN gates, respectively; the corresponding NAND
translations use 781 or 835 gates.  To avoid assuming a sign schedule in
elimination, \eqref{eq:axn-f19-factor} charges the larger value in each basis.
The generator and a separately implemented scalar evaluator both check every
one of the $19^3$ valid input triples.

For $\F_{19^4}$, the implementation uses the isomorphic sparse basis
\[
 \F_{19}[t]/(t^4-t-1).
\]
The element $13+13t+15t^2+2t^3$ is a root of the official defining
polynomial; its powers give the explicit basis change to $\F_{19}[S]$.
Two-level Karatsuba uses nine base-field multiplications, and $t^4=t+1$ makes
reduction additive.  The reported circuits for $c\mathbin{\pm}ab$ use 6,077
and 6,081 AXN gates, or 11,086 and 11,102 NAND gates.  Equation
\eqref{eq:axn-f194-factor} again charges the larger sign.  The reported
validation checks the complete scalar basis tensor, 4,096 random extension
products, the basis isomorphism, and both translated netlists.

Charging a signed multiply-accumulate for each multiplication upper-bounds
the local arithmetic primitive substituted into the published solver
schedules.  It does not upper-bound a complete solver once matrix
construction, control, memory, communication, and data movement are included.

\subsection{Same-basis AES calibration}\label{app:aes-calibration}

The reported implementation builds variable-key AES-128/192/256 one-pair
key-filter circuits using the NIST 113-gate S-box, the complete key schedule,
encryption, and a final ciphertext comparator.  To identify the actual key
rather than merely one matching a single pair, surviving candidates are
checked lazily on enough independent plaintext--ciphertext pairs to leave at
most $2^{-128}$ expected wrong survivors in the ideal-cipher model: two pairs
in total for AES-128 and three for AES-192 and AES-256.  Only survivors reach
the additional checks, so their expected work is negligible.  Every
first-pair circuit is validated on the FIPS test vector and eight additional
keys against a separately implemented AES evaluator.

\begin{table}[H]
\centering
\caption{Author-reported constructive AES calibration circuits.  Full charges
the entire key space for the first-pair filter; expected search is one bit
lower, and the lazy confirmation overhead is negligible under the stated
ideal-cipher model.  These circuits are not lower bounds on optimal AES
implementations and do not replace NIST's nominal references.}
\label{tab:aes-calibration}
\scriptsize
\setlength{\tabcolsep}{3.5pt}
\begin{tabular}{@{}c r r r r r r@{}}
\toprule
Level & pairs & NIST nominal & AXN gates/filter & AXN full & NAND gates/filter & NAND full\\
\midrule
I   & 2 & 143 & 30,081 & 142.88 & 106,616 & 144.70\\
III & 3 & 207 & 34,227 & 207.06 & 121,735 & 208.89\\
V   & 3 & 272 & 41,543 & 271.34 & 147,772 & 273.17\\
\bottomrule
\end{tabular}
\end{table}

The NIST-listed 28,600-gate AES-128 circuit has full-key-space exponent
$128+\log_2(28{,}600)=142.80$, numerically consistent with the nominal value
143.  The comparisons in the main text place arithmetic schedules beside
complete AES circuits and are therefore intentionally resource-asymmetric.

\section{Dimension Formulas}\label{app:dimensions}

For a Version~2.3 row, the following are parameter identities:
\[
 n=v+o,
 \qquad m_1=\left\lceil\frac{or}{\ell}\right\rceil,
 \qquad M=or\ell,
 \qquad N_0=n\ell,
 \qquad N_{\rm raw}=m_1\ell^2r^2.
\]
The number
$K=m_1\binom{\ell+1}{2}$ is the maximum formal quotient-feature count.
The check $\rank Q_R=K$ certifies that the public outer map is injective on
those formal quotient coordinates; it does not by itself prove that their
restrictions are independent polynomials.  For a structured slice, the
separate check $\rho_{\rm poly}(L)=K$ certifies its actual restricted
quadratic dimension.  For a fixed-skew system, any positive global polar-rank
lower bound implies independence of the homogeneous residual outputs; the
headline conditions are stronger still.  When
$\rank Q_R=K$ and $\rank C_{R,\mathbf v}=M-K$,
\[
 \dim\ker C_{R,\mathbf v}=N_0-M+K.
\]
If, in addition, $\rank\mathcal O_{R,\mathbf v}=m_1\ell^2$, then
\[
 \dim_{\F_q}H_{R,\mathbf v}=N_0-m_1\ell^2,
 \qquad
 \dim_A H_{R,\mathbf v}=n-m_1\ell.
\]
When the selected $A$-stable special-XL slice also passes
$\rho_{\rm poly}(L)=K$, it has $A$-dimension $s$, base-field dimension
$h=\ell s$, and retry exponent $\tau=K-h$.

\begin{table}[tbp]
\centering
\caption{Structured-lift dimensions for the six Version~2.3 $\ell=4$ rows,
conditional on the rank preflights stated above.}
\label{tab:dimensions}
\small
\setlength{\tabcolsep}{5pt}
\begin{tabular}{@{}l r r r r r@{}}
\toprule
Parameters & $M$ & $K$ & $\dim_A H$ & $s$ & $\tau$\\
\midrule
$(28,5,4,4)$ & 80  & 50 & 13 & 10 & 10\\
$(28,4,4,5)$ & 80  & 50 & 12 & 10 & 10\\
$(40,7,4,4)$ & 112 & 70 & 19 & 15 & 10\\
$(38,5,4,5)$ & 100 & 70 & 15 & 15 & 10\\
$(50,9,4,4)$ & 144 & 90 & 23 & 19 & 14\\
$(52,6,4,6)$ & 144 & 90 & 22 & 19 & 14\\
\bottomrule
\end{tabular}
\end{table}

\section{Cost-Search Certificate}\label{app:cost-certificates}

For each $m_1\in\{5,7,9\}$ and every admissible integer $s$, the reported
certificate search
expands \eqref{eq:hilbert-series} exactly.  Here admissible means
$1\le s\le\min\{\dim_A H,\lfloor K/4\rfloor\}$.  The $K/4$ bounds are
$12,17,22$.  They bind or tie for both $m_1=5$ rows and both $m_1=9$ rows.
For $m_1=7$, the bounds are $17$ for $(40,7,4,4)$ and $15$ for
$(38,5,4,5)$, the latter coming from $\dim_A H$.  The search reaches
the larger group bound and then restricts the certificate separately for each
parameter set; both $m_1=7$ rows attain the same optimum at $s=15$.
The monomial count
$M_s(\mathbf d)$ is coordinatewise increasing, so any multidegree that could
match or beat the incumbent lies in a finite rectangular box.  By symmetry it
suffices to enumerate sorted multidegrees.  The search checks 3,690, 11,217,
and 27,449 candidates for $m_1=5,7,9$, respectively.
 
The optima, unique up
to block permutation, are
those in \Cref{tab:xl-optima}.  No zero signed coefficient occurs anywhere
in these bounded searches, so using a strictly negative rather than a
non-positive trigger does not change the result.  The optimal coefficients are
\[
 -14{,}677{,}784,\qquad
 -124{,}983{,}335{,}190,\qquad
 -1{,}164{,}044{,}975{,}088{,}360.
\]
The corresponding Macaulay column counts are
\[
 23{,}417{,}049{,}656,
 \quad 44{,}237{,}557{,}981{,}725{,}696,
 \quad 236{,}094{,}291{,}515{,}544{,}000{,}000.
\]
Substitution of the full-rank retry bound
$\mathsf{Retry}_{K,\tau}$ into \eqref{eq:special-xl-cost} gives published-model exponents
128.81713, 171.68562, and 214.12301.  Replacing $b_4$ by the
author-reported signed AXN circuit
$g_4^{\mathrm{AXN}}$ gives converted arithmetic exponents 132.17187,
175.04035, and 217.47775.  The NAND sensitivity gives 133.04031, 175.90879,
and 218.34618.

\section{Pinned Generic-MQ Objective}\label{app:generic-mq-certificate}

For completeness, this section records the finite optimization used for every
generic-MQ row.  This removes any dependence on an implicit software default
or an unquoted ``Appendix A convention.''  We first factor out the common
cost assigned to one field operation.  For $n,m\ge1$, define the arithmetic
schedule
\[
 \widetilde{\operatorname{CFXL}}_q(n,m)
 =\min_{0\le k'<n}
 q^{k'}\,3
 \binom{n-k'+D}{D}^{\!2}
 \binom{n-k'+2}{2},
\]
where, for each candidate $k'$, $D$ is the least nonnegative integer such
that
\[
 [t^D]\frac{(1-t^2)^m}{(1-t)^{\,n-k'+1}}\le0.
\]
A candidate for which no such $D$ occurs is omitted.  Hashimoto's
underdetermined-system objective is
\[
\begin{split}
 \widetilde{\mathcal H}_q(n,m)=\min_{a,k}\ \{&
 (m-a-k+1)\widetilde{\operatorname{CFXL}}_q(a,a)\\
 &+q^k\bigl(
 \widetilde{\operatorname{CFXL}}_q(a-1,a-1)
 +\widetilde{\operatorname{CFXL}}_q(m-a-k,m-a)\bigr)\},
\end{split}
\]
over integers $1\le a\le m$, $0\le k\le m-a$, subject to
\[
 n\ge(a+1)(m-k-a+1),
 \qquad
 n\ge a(m-k)-(a-1)^2+k.
\]
The minimization is further restricted to points for which the third CFXL
call is defined: either $(a,k)=(m,0)$, giving
$\widetilde{\operatorname{CFXL}}_q(0,0)$, or $a<m$ and
$0\le k<m-a$.  Thus no value
$\widetilde{\operatorname{CFXL}}_q(0,j)$ with $j>0$ is implicitly used.
The boundary conventions are
\[
 \widetilde{\operatorname{CFXL}}_q(0,0)=1,
 \qquad
 \widetilde{\operatorname{CFXL}}_q(1,1)=q,
\]
where the second convention assigns one primitive unit to direct evaluation
at each field element because the displayed Hilbert trigger does not occur.
Within a CFXL call, an
exact tie is resolved by the smallest $k'$; an exact outer tie is resolved
lexicographically by $(a,k)$.  The search compares exact integer arithmetic
schedules and rounds only the final base-2 logarithm.  The Model schedule is
$b_1\widetilde{\mathcal H}_q(n,m)$; the AXN and NAND schedules replace $b_1$
by $g_1^{\rm AXN}$ and $g_1^{\rm NAND}$, respectively.
For $q=19$, the second boundary convention is the elementary
$\operatorname{MQ}(19,1,1)$ solver that exhaustively evaluates all 19 field
elements.

\begin{table}[H]
\centering
\caption{Exact witnesses used by the pinned generic-MQ search.  Here $k'$ is
the inner specialization minimizing the third call
$\widetilde{\operatorname{CFXL}}_q(m-a-k,m-a)$.
The first three rows are the generic symmetry-quotient baselines used for the
structured rows' internal estimator gains; grouped residual systems share the
same optimum.  Exact-integer comparison is used to choose among points that
agree to all displayed decimal places.}
\label{tab:generic-mq-witnesses}
\scriptsize
\setlength{\tabcolsep}{3.5pt}
\begin{tabular}{@{}l r r r r@{}}
\toprule
Residual system & $(a,k,k')$ & $D$ & Model exponent & AXN arithmetic schedule\\
\midrule
$50/102,\ 50/98$ structured baseline & $(1,1,10)$ & 12 & 136.039949042 & 140.34608\\
$70/146,\ 70/142$ structured baseline & $(1,1,11)$ & 18 & 180.595403793 & 184.90153\\
$90/182,\ 90/178$ structured baseline & $(1,1,18)$ & 19 & 224.687086566 & 228.99321\\
\midrule
$48/112$ fixed-skew  & $(2,10,1)$ & 11 & 129.440299085 & 133.74643\\
$72/168$ fixed-skew  & $(2,15,0)$ & 16 & 183.438029550 & 187.74416\\
$96/256$ zero-offset & $(2,10,6)$ & 23 & 235.194737432 & 239.50086\\
$96/224$ fixed-skew  & $(1,1,18)$  & 21 & 237.843096284 & 242.14922\\
\bottomrule
\end{tabular}
\end{table}

\section{Same-Basis Ordered-Label Estimator Baseline}\label{app:ordered-baseline}

The comparison with the official Version~2.3 table mixes cost conventions, so
we also recost a reproducible \emph{estimator baseline} under exactly the same
AXN arithmetic basis as the selected conditional schedules.  This baseline feeds the full
ordered-label equation/variable dimensions from the earlier affine-column
accounting into the pinned generic-MQ optimizer and replaces each charged
base-field operation by $g_1^{\mathrm{AXN}}=441$.  An attacker can indeed feed
the unreduced affine-column equations to a generic solver.  What is
counterfactual is the estimator's treatment of the ordered labels as though
they represented independent generic quadratic features: on every symmetric
key, reversed labels define the same polynomial by
\Cref{thm:symmetric-square}.  Thus the table is a controlled historical
estimator comparison, not evidence for a second quadratic space of the stated
ordered-label rank.  It reports every row rather than only the range quoted in
the introduction.

\begin{table}[H]
\centering
\caption{Historical ordered-label generic-MQ estimator and selected
conditional AXN arithmetic schedule, both recosted in the same AXN arithmetic
basis.  ``Estimator input'' gives the
equation/variable dimensions of the unreduced affine-column equations.  The
system is executable, but treating its reversed labels as independent generic
quadratic features is counterfactual on a symmetric key.}
\label{tab:ordered-baseline}
\footnotesize
\setlength{\tabcolsep}{4.1pt}
\renewcommand{\arraystretch}{1.08}
\begin{tabular}{@{}c l c r r r@{}}
\toprule
Level & Parameters & \shortstack{Estimator\\input} & \shortstack{Ordered-label\\AXN} &
\shortstack{Selected cond.\\AXN schedule} & \shortstack{Estimator\\difference}\\
\midrule
I   & $(28,5,4,4)$  & $80/132$  & 206.69 & 132.17 &  74.52\\
I   & $(48,16,2,2)$ & $64/128$  & 171.77 & 133.75 &  38.02\\
I   & $(28,4,4,5)$  & $80/128$  & 206.69 & 132.17 &  74.52\\
\midrule
III & $(40,7,4,4)$  & $112/188$ & 277.52 & 175.04 & 102.48\\
III & $(72,24,2,2)$ & $96/192$  & 242.15 & 187.74 &  54.41\\
III & $(38,5,4,5)$  & $100/172$ & 250.29 & 175.04 &  75.25\\
\midrule
V   & $(50,9,4,4)$  & $144/236$ & 345.49 & 217.48 & 128.01\\
V   & $(96,32,2,2)$ & $128/256$ & 311.08 & 242.15 &  68.94\\
V   & $(52,6,4,6)$  & $144/232$ & 347.13 & 217.48 & 129.66\\
\bottomrule
\end{tabular}
\end{table}

The estimator-difference endpoints, rounded to five decimals, are
$38.02124$--$129.65539$ bits.  This comparison
still inherits the generic-MQ and special-XL solving models, but it removes the
field-operation and gate-basis mismatch from the comparison itself.

\section{Reproducibility Contract}\label{app:artifact}

Machine-readable certificates for the reported computations exist, but this
manuscript does not provide an immutable archive locator for the complete
supporting bundle.  The computational claims are therefore labeled
author-reported and are not used as proofs of a global rank condition.  A
complete archival package enabling independent reproduction should provide
all of the following:
\begin{itemize}
\item the exact paper source and bibliography, an immutable repository commit
or archive identifier, and a root manifest hashing every released file;
\item the public-key generator and unmodified verifier, the official
known-answer inputs and outputs, and the literal square and rectangular
output-order regression tests;
\item implementations of the quotient, affine feature map, structured offsets,
stable kernel, and block-symmetry rejection test, with every public choice, seed, key
identifier, slice basis, and preflight transcript needed to reproduce the
reported 121 structured instances, 13 slices, and 60 fixed-skew instances,
as well as the complete 361-relation $\ell=2$ quotient-rank transcript;
\item every direction vector and exact rank transcript underlying the reported
cross-polar, directional-derivative, fixed-skew, and zero-offset campaigns,
together with the field representation, matrix dimensions, storage format,
elimination algorithm, pivot convention, and a separately implemented rank
cross-check;
\item the existing cost-search code and exact-integer certificates for
\Cref{app:cost-certificates,app:generic-mq-certificate}, including tie
breaking, endpoint handling, table-generation commands, and machine-readable
outputs, together with the exact path and revision of every historical
analysis script used for a floor-versus-ceiling comparison;
\item the rejection-sampling and target-generation code, the per-parameter
chunk layouts and acceptance probabilities, the duplicate-digest and
chosen-prefix schedule, and the exact Level-V salt-space and prefix-budget
calculation in \Cref{app:zero-offset};
\item for each Macaulay experiment, the complete polynomial profile, selected
multidegree, monomial order, row and column counts, number of nonzeros, field
encoding, matrix-construction and storage algorithms, exact-rank or
kernel algorithm, random seed, software versions, invocation command,
hardware description, the full kernel-dimension histogram, and the
base-field root count or other record certifying $Z=1$ for each planted case
used to assess $p_{\rm ker}^{(1)}$;
\item source or netlists for every field and AES circuit, circuit conventions,
known-answer inputs, separately implemented evaluators, validation
transcripts, and hashes; and
\item environment and dependency lockfiles plus a one-command workflow that
rebuilds every paper table from the archived raw records.
\end{itemize}
The intended inventory spans the computational claims throughout this
manuscript: exact public preflights and finite-field ranks over $\F_{19}$,
bounded cost searches, smaller Macaulay controls, fully enumerated small-field
probability checks, target-generation calculations, and circuit validation.
Until the certificates and supporting records are attached to an immutable
locator cited here, those computations remain author-reported and are not
independently reproducible from this manuscript alone.

\begingroup
\footnotesize
\printbibliography
\endgroup
\end{document}
